\chapquote{Coronations}{Therefore I am sure that this, my Coronation, is not the symbol of a power and a splendour that are gone but a declaration of our hopes for the future, and for the years I may, by God's Grace and Mercy, be given to reign and serve you as your Queen. \\ \textemdash Queen Elizabeth II}
\phantomsection
\label{cha:Coronations}

\begin{multicols}{2}\raggedcolumns

\section*{Creating a Princess} 
\phantomsection
\label{sec:CreatingAPrincess}

\subsection*{Step One: The Story Seed} 

A story seed is an idea for a game in Princess: the Hopeful, and it's the first step to play the game. The story seed can either be the Storyteller's pitch or created collaboratively by the group. Either way, it's a good idea to have a quick outline of what kind of game you're playing before creating characters. 
%Editor: Italicize line?

Is your game about inspiring your classmates in a demoralized school? Then everyone has to be a student who cares about education and has a background which places them in a poorly performing school. Is your game about a Nakama hand-picked by a Queen to track down and destroy a powerful Darkspawn hiding in the jungle? Then everyone should create a character who has useful skills for the quest and nothing in their backstory which will prevent them from going away for a week.

The story seed doesn't have to be complex or well-defined. It just needs to be enough for players to design their character's background and their character's abilities appropriately. 

Some story seeds result in narrower mechanical choices. A story where the characters are trying to turn around a failing inner-city school has a lot of flexibility. Characters could focus on helping their classmates study, on working with the faculty and local government, or even on fixing the building themselves. However, with a story seed about fighting Darkspawn, everyone will probably want at least one offensive Charm and the defensive Charms Celestial Grace and Barrier Jacket. With upgrades, that's a serious investment of a starting Princess's dots in Charms. It would not be odd or inappropriate if such a story seed resulted in a group made up mostly, or even entirely, of Champions. When this happens, the Storyteller should consider giving additional Charm and Merit dots to the players so they can cover the basics add still have room to give their character's unique talents.

\subsection*{Step Two: Character Concept}

The second step is to try and distill your character down into a single concept. You want something that you can easily hold in your head. This concept will serve as the anchoring center to your character. When you start adding details later, it should be quick and straightforward to ask yourself \textquotedblleft does this detail strengthen my character concept?\textquotedblright 

A good character concept focuses on who your character is. Is your character's self-identity defined more by her mundane side or her supernatural side? Or does she balance both? Because Princesses invest so much of themselves in their work, it is often hard to separate who a Princess is from what she does, but try to cut out the specifics and drill down to the core of your concept. Consider playing a young politician instead of the student president.  

Characters often change during character creation as ideas bounce around. Once play starts, circumstances can change fast. It's good to have a little flexibility in your core concept. 

\subsection*{Step Three: Fleshing It Out}

Are you playing a Prince or a Princess? Does your character fight evil in a frilly sailor fuku or spandex and a cape? 

Now that you have your core character concept, you can add everything else: appearance, personality, background, family, friendships, and perhaps romances. In short, it's time to write your character as a character rather than a collection of attributes and a list of magical powers. 

Some players can write for pages while others will prefer a few paragraphs. There's no right or wrong answer, though that it never hurts to discuss things with your group and make everyone's style fit together.  


\begin{mybox}


\subsubsection*{Burdens of the Crown}

There are several persistent conditions that are appropriate to take at character creation. Players may wish to consider 
\begin(itemize)
\item Adult Supervision
\item Body Dysmorphia
\item Embarrassing Phylactery
\item Heartfelt Vow
\item Name of the Moon
\item Truly Mundane
\item Sensory Overload
\item Slow Transformation
\item Still Phylactery
\item Supernatural Heraldry
\item Two Hearts, One Soul
\item Unstable Transformation

The usual persistent conditions from the Chronicle of Darkness book such as Addicted, Blind, Crippled or Mute, are of course appropriate options for Princesses and function normally. At the player's discretion, a condition may only affect the character in one of her two forms, so long as it makes sense.
%Editor: The list of conditions was really long, so I figured it would be worth putting it in a bulleted list. No idea how it'll look in the text once compiled, though, so feel free to change it back. 

\end{mybox}


\subsection*{Step Four: Virtues, Vice, Aspirations, and Dreams}   

A Princess is a magical being, but she's still human. She has human wants and human needs. To complete your character, assign two — yes, two — Virtues and one Vice. It is recommended that Princesses do not take the \textit{Hopeful} Virtue. Nearly every Radiant Princess demonstrates that Virtue every day of their lives. Players should use Virtues as an opportunity to make their characters stand out from the stereotypical Noble.
%Editor: Note to self: Check if italicize Virtues? 

You should also think about what your character wishes to accomplish, which means it's time to choose her Aspirations and Dreams. Aspirations (see \textit{Chronicle of Darkness} page 28) are a Princess's goals, while Dreams (see page \pageref{subsec:Dreams}) are variant Aspirations that provide an additional benefit when a Princess's goal is to help others through her Calling. A Princess has a total of five Aspirations and Dreams, two of which must be Dreams; the remaining three can be either. 

As a guideline, you should try to have at least one Aspiration or Dream that reflects your character's self-image, and you should have at least one Aspiration or Dream that ties you directly into the story seed. 

Using the examples above, a player in the game about turning around a demoralized school could consider a Dream such as \textquotedblleft help Harriet realize she's better than her mother says \textquotedblright. In one sentence, the player has chosen an ambition which fits the story's core themes, connected their character to two NPCs (one of whom is an antagonist), and identified social combat as a key skill.

The example of a quest into the jungle might be too specific to easily create distinctive Aspirations, and there's nothing wrong with that. 

\subsection*{Step Five: Add Attributes}

Now is the time to define your character mechanically, starting with the most fundamental aspects. First, choose which of the three categories (Mental, Physical, and Social) will be your character's primary, secondary, and tertiary focus. Then place five dots (in addition to your starting dot in every attribute) into your primary category, split between the three Attributes as you wish. Place four dots into your secondary focus and three into your tertiary focus. 
%Editor: Need to confirm whether to cap all the things in here. Look into COD book? Confirmed: always cap Attributes 

When choosing your Attributes, ask yourself which will be most helpful for your character's Calling. Then, decide if she's good at them or if she relies on her magic. Starting as a competent and successful Noble and starting with lots of room to grow are both valid choices. 


\bigsidebarstart

\subsubsection*{Daddy's Little Princess}

The mahou shoujo genre has a long and proud tradition of child and early teen protagonists. If you wish, you may combine \textit{Princess} with \textit{The World of Darkness: Innocents}. The \textit{Princess} template will function as intended if applied to a character built using the rules in \textit{Innocents}. Any Charm that uses Science or Academics will use the Study skill instead. (If you prefer, you can instead split your character's Study dots into Academics and Science upon Transformation). \textit{Innocents} gives characters different starting skill dots depending on their age (see \textit{Innocents}, page 47). However, \textit{Innocents} also suggests that Storytellers might prefer to start all the player characters on an equal footing. When playing \textit{Princess}, Storytellers can give extra Transformed skill dots (which do not count toward the maximum) as an additional way of achieving balance. 

Child characters are limited to three dots in Skills unless they have one of the three Prodigy Merits. This limitation does not apply to Transformed Dots, and the cap on mundane Skills is increased by one at Inner Light four and five. Also, so long as a Princess has at least one relevant Transformed Dot, she is considered an adult for rules related to that Attribute or Skill while she is Transformed. The same applies to anyone who has 6+ dots in an Attribute or Skill, including anyone under the influence of an appropriate Bless Charm. 

For example, think of a Princess with a Transformed dot in Strength. 
\begin(itemize)
\item She would not subtract an adult target's Stamina from attack rolls.
\item She would use the adult strength chart.
\item In a Contested roll against an adult, the adult would not gain 8-again if Strength is part of the Princess's dice pool.

A Princess's magical body is more than a match for size and leverage. (A complete list of special rules for adult-child interactions can be found in \textit{Innocents}, page 193). How this manifests is up to the player. A Princess with Transformed Composure might acquire adult emotional maturity every time she Transforms or overcome her fears through a supernaturally potent child-like sense of adventure. 

With a little cunning, Princesses can access Social Merits that are usually off-limits to children. A Troubadour could earn Resources by working as an artist for commission. Nobody can see that she's nine years old over e-mail, and her work could put adult professionals to shame. Similar tricks and the occasional Appear Charm to look like an adult could allow a Princess to acquire 4+ dots in Allies or Contacts. 

All this assumes that the young Princess doesn't become an adult upon Transformation. If she does, then just treat the Transformed Princess like any other grown-up.

The other characters' reaction to a Noble child should be negotiated between the group to identify the correct trade-off between realism and fun. Most adults, Noble or otherwise, would not approve of children running toward danger, but a Princess sitting at home and waiting until she's 18 will not create a dramatic story. A child Princess could be treated equally to the adults for no reason except that the character wouldn't be fun to play otherwise. Alternatively, the child protagonist being sent home only to sneak out again and save the adults is a time-honoured trope. Likewise, a mortal (or Sworn) who's constantly dragged into his Noble kid sister's dangerous magical adventures because he's powerless to stop her could make an interesting character, especially if the rest of the players don't acquire magic that overshadows his special power of being a grown-up. 
\bigsidebarend

\subsection*{Step Six: Add Skills}

Skills are chosen like Attributes. From the Mental, Physical, and Social Skills categories, pick which will be your primary, secondary, and tertiary focus. You get 11 primary dots, seven secondary dots, and four tertiary dots to spend. Unlike Attributes, the first dot in a skill is not free. 
%editor: Check if focus wording is right here. 

As with Attributes, it's worth thinking about whether you wish to start off competent or if you wish for your character to begin the game dependant on her magic. 

\subsection*{Step Seven: Add Skill Specialities}
 
Skills are extremely broad categories, which means your choice of Skills is an imprecise reflection of your character. Does three dots in Expression make her a singer or a poet? Specialities let you correct this. Each speciality is a one- or two-word description that can expresses what fields your character specializes in. Pick three specialities for your character. 
 
There are many simple Charms available to starting characters that can enhance a Princess's specialities. It might be worth examining them before selecting your specialities.

\subsection*{Step Eight: Add a Calling and a Transformed Attribute Dot}

The eighth step in creating a Princess is to choose her Calling. A Calling is the way a Princess brings hope to the world. Her Calling determines which Oaths are part of her Belief, what duties she performs to regain Wisps, which Attributes are increased when she Transforms, and which Charm families she has an affinity for.
%Editor: Changed "first step" to "eighth step" and removed part of the last sentence for brevity. 

\begin{itemize}
\item Champion: Defenders of the weak and helpless. (Strength or Resolve)
\item Grace: Elegant and noble, the voices and diplomats of the Reborn. (Presence or Composure) 
\item Mender: Those who take wounded flesh and broken hearts and make them whole. (Intelligence or Dexterity)
\item Seeker: Hunters of truth, banishers of lies, revealers of the hidden. (Intelligence or Wits)
\item Troubadour: Artists and exemplars who inspire greatness in others. (Presence or Dexterity)
\end{itemize}

\subsection*{Step Nine: Add a Queen}

Once, the five Radiant Queens governed their Kingdoms in the name of the Light. Now, they are dead, but their minds and souls live on in the Dreamlands. From beyond the veil of sleep, the Queens mentor the Radiant. A Princess can accept a position in a Queen's court, a bond of both personal loyalty and mystical affinity that enhances a Princess's magic. 
%Editor: Is "veil of sleep" a proper noun, or is it just a fancy way of saying the Queens are dead? Figured it was the former, but wanted to check. 

There are five Radiant queens that a player character can follow. Rarely, some Princesses choose to not join any Court. There are also three Twilight Queens, who are not usually available to player characters. 

\begin{itemize}
\item Clubs: The Queen of Wilds leads the coven of Princesses who seek to live in harmony with the world.
\item Diamonds: The Queen of Lights leads the conclave of progressive-minded Princesses who believe in constantly adapting and refining the kingdom's techniques.
\item Hearts: The Queen among Queens leads the court of Princesses who focus on communities and their place within society.
\item Spades: The Queen of Knaves leads a confederacy of winking scoundrels and good-hearted rogues who believe the best thing to do with rules is break them.
\item Swords: The Queen of Heroines leads the company of Princesses who believe that inner strength and passion should always light the way. 
\end{itemize}

\bigsidebarstart
 
\subsubsection*{Choosing a Queen}
Princesses rarely agree with her Queen on every single point, and almost all Princesses have some affinity for the teachings of several Queens. This can make it hard for Princesses (and players) to know which Court is right for them. Should a naturalist join the Court of Diamonds with her fellow women of science or the Court of Clubs with her fellow nature-lovers? Such a Princess will probably show her shared ideals by acquiring both the Acqua and Legno Invocations, but the player still has to decide which Queen goes on the top of the character sheet. How would the player figure that out? They can decide which Queen most closely shares the Princess's approach to make a brighter tomorrow:

\begin{itemize}
\item If you believe that self-discovery and mutual understanding bring people together in harmony, you will always be welcome in the Court of Clubs.
\item If you believe that clear and careful thought will lead to a better tomorrow, you'll find like minds in the Court of Diamonds.
\item If you believe that a Princess's role is to build and lead communities into the Light, then you have a place in the Court of Hearts.
\item If you believe that the world's ills are best solved outside of any boxes, then you'll fit right in at the Court of Spades.
\item If you believe that the world needs individual heroes led by all-consuming passion, then you can proudly join the Court of Swords. \\
\end{itemize}

The Twilight Queens also have an approach:

\begin{itemize}
\item If you believe that destroying the Darkness matters more than helping the world, you are ready to enlist in the armies of Storms.
\item If you believe that being a Princess makes you the most important and special person in the world, and a utopia is just going to happen because you're around, then consider the Queen of Mirrors.
\item Finally, the Queen of Tears would say that if you believe in protecting your people no matter the cost, bring them to Alhambra where they will be safe. Remember: of all the Queens, only the Lady of Alhambra's kingdom survives to this day. She asks, and offers, more than any other Queen. A Princess who believes in protecting her people no matter what should still think if she wishes to be a part of Alhambra. If not, she should choose another Queen or perhaps become a Courtless Lacrima specialist. 
\end{itemize}



\bigsidebarendw

\subsection*{Step Ten: Add Invocations} 

All Princesses draw magic from their Beliefs and hopes, but when a Princess becomes a Queen, her Beliefs echo across the souls of humanity. Any Princess can feel that echo within her soul, and she can draw on it to enhance her own magic. Each Princess starts with three dots in Invocations, one of which must be placed into her Queen's favored Invocation. Courtless Princesses do not have this restriction, but they may not place all three of their dots in the same Invocation. 

There are eight established Invocations, each favored by a Queen.

\begin{itemize}
\item The Queen of Clubs teaches the principles of harmony and inner strength in the Invocation of Legno.
\item The Queen of Diamonds focuses the lights of knowledge and reason into the Invocation of Acqua.
\item The Queen of Hearts builds her society upon the foundation of the Invocation of Terra.
\item The Queen of Spades teaches the Invocation of Aria, the enviable lightness of being.
\item The Queen of Swords raises passionate devotion into magic in the Invocation of Fuoco.
\item The Queen of Tears has embodied the laws and requirements of her city within the Invocation of Lacrima.
\item The Queen of Storms distills rage against Darkness into the Invocation of Tempesta.
\item The Queen of Mirrors uses the Invocation of Specchio to cultivate those who may be the Kingdom’s heir.
\end{itemize}

\subsection*{Step Eleven: Add Charms and Transformed Skills}

In her perfected self, every Princess has exceptional skill and can use Charms and special talents. A Princess starts with three Dots in Transformed Skills and nine dots in Charms. At least three dots must be for Charms in the three families the character's Calling has affinity for; the others may come from any family. You may take upgrades for any Charm you know; each upgrade costs 1 of your starting Charm dots. If you take an invoked Charm or upgrade, you must have the prerequisite dots in the Invocation. The ten families of Charms are

\begin{itemize}
\item Appear Charms change a Princess’s appearance and produce other illusions. (Affinity for Seekers and Troubadours)
\item Bless Charms enhance people, pushing their abilities to greater levels. (Affinity for Champions and Graces)
\item Connect Charms trace and manipulate interpersonal connections. (Affinity for Graces)
\item Fight Charms enhance a Princess’s fighting ability. (Affinity for Champions)
\item Govern Charms control supernatural powers and those who possess them. (Affinity for Graces and Seekers)
\item Inspire Charms lift people’s hearts, spark their passions, and exhort them to action. (Affinity for Troubadours)
\item Learn Charms confer knowledge and reveal what has been hidden. (Affinity for Seekers)
\item Perfect Charms grant new personal abilities. (Affinity for Champions and Menders)
\item Restore Charms repair damaged objects, heal the injured, and tend to the sick. (Affinity for Menders)
\item Shape Charms shape materials and make things out of nothing. (Affinity for Menders and Troubadours)
\end{itemize}

To take a 2-dot Charm, you must also take a different 1-dot Charm of the same family. To take a 3-dot Charm, you must take 3 dots worth of Charms in the same family. (For this purpose, an upgrade counts as a single dot.) If the Charm has a prerequisite on an Invocation, you may treat it as part of a pseudo-family and use it to unlock access to other Charms that depend on the same Invocation. 
%Editor: Check wording on "prereq on an Invocation," seems weird? 

If you wish, you may also trade Charm dots for additional Transformed Attributes and Skills. You may not spend your 3 affinity Charm dots this way. A Charm dot is worth 1 Transformed Skill dot, and three Charm dots are worth one Transformed Attribute dot.

\begin{mybox}


\subsubsection*{Free Affinity \textemdash\, an optional rule}
Some players have character concepts that don’t fit neatly into the five Callings and need Charms from two families favored by no single Calling. To facilitate such characters, the Storyteller may rule to play each Calling with two fixed Charm affinities and a free pick for the third affinity. Once chosen, the third affinity is fixed and can no more be changed than a Princess's Calling.

Under this scheme, the affinities are

\begin{itemize}
\item Champion: Fight and Perfect
\item Grace: Connect and Bless
\item Mender: Restore and Shape
\item Seeker: Learn and Govern
\item Troubadour Inspire and Appear
\end{itemize}

\end{mybox}


\subsection*{Step Twelve: Add Inner Light}

A Princess's magic shines out of her soul to illuminate the world. The raw power of her hopes and feelings is represented by her Inner Light. Princesses start with one dot in Inner Light, but she may spend five Merit dots to increase it to two, or 10 Merit dots to increase it to three. 

\subsection*{Step Thirteen: Add Wisps}

Wisps are the fuel for a Princess, a measure of both raw magic and emotional energy. A Princess begins with half her starting Wisp pool, as measured by Inner Light, plus the higher of her Circle and Mandate.

\subsection*{Step Fourteen: Add Belief} 

Belief is the Integrity trait for Princesses. It measures her confidence that the world can be improved and how moral and principled her actions are: a Princess is innately moral, but as her confidence breaks, she becomes desperate enough to consider unprincipled or outright immoral solutions to her problems. Starting characters begin with a Belief of 7, representing idealism and confidence with some concessions to practicality. 

During character creation, ask yourself the following questions. These questions will define a Princess's belief system. These questions are not intended to identify specific compromises, though you will probably discover some. Rather, they are to help you create a moral framework that can be used to identify new compromises and which can be modified as dots of Belief are lost or gained. 

\begin{itemize}
 \item \emph{What does your character's Calling mean to her?} Is your Champion a modern day knight in shining armour who uplifts and inspires the public with her virtuous deeds, or does she fight a thankless secret war in the shadows so civilians can sleep soundly? Is your Grace a mentor, a judge, or a leader? A character's core concept can be the starting point of her Beliefs. Her ideals and principles explain why she chose this mission instead of another. 
 \item \emph{What does your character believe is wrong with the world?} Princesses are trying to improve the world, so it stands to reason that the Hopeful see the world as needing to be fixed. Why is the world in such a bad state? Is it all down to public ignorance and apathy, or is it because humanity lost touch with nature? No matter how justified the means, any Princess who uses or encourages what she believes to be the world's flaws is betraying her principles. 
 \item \emph{What is your character's perfect world?} Gandhi apocryphally said \textquotedblleft be the change you want to see in the world.\textquotedblright What is your character's ideal world? Is it a vegan nature utopia, or something more along the lines of an technological post-scarcity society? How would a person from such a society act? What unusual and idealistic cultural values does your Princess follow, and how does her utopia give context to her actions?
 \item \emph{Who influenced your character's philosophy?} Did your character learn her idealism from her parents? Her wise and outrageous grandmother? A teacher? Think a little about their relationship and her mentor's personality. How do they influence your character's relationship to her ideals? If your character grew up without guidance, it is perfectly acceptable to use someone from a past life. 
 %Editor: Is "outrageous" meant to be "courageous"? Just seemed like the two adjectives of the grandma don't match. 
 \item \emph{What is the best argument against her Beliefs she has ever heard?} Ultimately, utopia is subjective. A life of peace and tranquillity would bore some people to tears. Every Princess is going to find people who disagree with her vision, from practical problems to ideological conflicts. Of all the people who disagreed with her, who left your character speechless with a perfect argument or irrefutable fact, and how did this affect your character?
 \item \emph{What does your character care about more than her Belief?} Princesses are still people, and people's moral codes rarely survive when rigorously applied to every situation. Would your character do anything to save her boyfriend, or perhaps herself? Most people have an instinctive wish to survive. The goal here isn't to create a hypocrite, but to give depth to a character. 
\end{itemize}

\noindent A Princess risks losing Belief when her actions go against her own principles. Now that you have fleshed out your character's belief system, it should be straightforward to identify ways she might violate it. Further details on the mechanics for losing dots of Belief can be found on page \pageref{subsec:Belief}.

\subsection*{Step Fifteen: Add Merits} 

As a supernatural being, a Princess is given three merit dots to spend at character creation, for a total of ten merit dots. Her magic makes it easier to acquire resources and skills, while her shining heart wins friends and allies. Additional Merits unique to the Hopeful are presented in Chapter Three. Additionally, some Merits presented in \textit{God-Machine Chronicle} are inappropriate or work differently for the Hopeful. This is discussed in more detail in Chapter Three. 
%Editor: "God Machine Chronicle" is talking about this book, right? https://archive.org/details/god-machine-chronicle/Chronicles%20of%20Darkness%20-%20The%20God-Machine%20Chronicle

%\subsection*{Step fourteen: Add Finishing Touches} 

%A Princess is a magical being, but she's still human. She has human wants and human needs. To complete your character assign two, yes two, Virtues and a Vice. It is recommended that Princesses do not take the \textit{Hopeful} Virtue. Nearly every Radiant Princess demonstrates that Virtue every day of their lives, players should use Virtues as an opportunity to make their characters stand out from the stereotypical Noble and avoid taking a trait which would be redundant on top of Belief.

%As befitting a group referred to as the Hopeful Princesses have strong Aspirations, as well as Dreams which are described later in this Chapter. A Princess begins with three Aspirations and two Dreams, the player may trade Aspirations for additional Dreams at a one for one ratio. 
%Editor: Skipping this commented-out section. 

\bigsidebarstart

\subsection*{Character Advancement}

\begin{quote}There is nothing noble in being superior to your fellow man; true nobility is being superior to your former self.\\
\rightline{\textemdash\, Ernest Hemingway} \end{quote}
%Editor: "Hemingway" is longer than the gray box on the page. Anything that can be done about that, formatting-wise? 
\phantomsection
\label{sec:CharacterAdvancement}

    \centering
    \begin{tabular}{|l|l|}
        \hline
        Mundane Attribute       & 4 Experience per dot                                \\ \hline
        Mundane Skill           & 2 Experience per dot                                \\ \hline
        Transformed Attribute   & 3 Experience per dot                                \\ \hline
        Transformed Skill       & 1 Experience per dot                                \\ \hline
        Transformed to Mundane  & 1 Experience per dot  \\ \hline
        Affinity Charm          & 1 Experience per dot                                            \\ \hline
        Other Charm             & 1 Experience + 1 Experience per dot                                           \\ \hline
        Upgrade          	 & 1 Experience per Upgrade                                                   \\ \hline
        Favored Invocation     & 2 Experience per dot                                        \\ \hline
        Other Invocation        & 3 Experience per dot                                        \\ \hline
        Merit                   & 1 Experience per dot                                        \\ \hline
        Inner Light             & 5 Experience per dot                                       \\ \hline
        Belief                  & 2 Experience per dot                                    \\ \hline
        Willpower               & 1 Experience per dot (bought back after sacrifice)                     \\
        \hline
    \end{tabular}	
    
    \bigsidebarend

\section*{Traits} 
\phantomsection
\label{sec:Traits}

Becoming a Princess does not just grant you magical abilities, it is a change to your fundamental being. Many Princesses use the phrase \textquotedblleft magical girl\textquotedblright\, to emphasize this point. A Princess is not a girl who uses magic, she \textit{is} magical. Magic is as much a part of her as blood is a part of any human. As a Princess grows in wisdom and experience, she learns to use her nature to its fullest extent, but she also must face its drawbacks. 

\subsection*{Inner Light}
\phantomsection
\label{subsec:InnerLight}
When a Princess's power first blossoms, it is barely the tiniest of sparks deep within her soul. However, as time passes and the ancient spirit within her awakens, her power begins to shine like a roaring fire, strong enough to support great magics and bright enough to become an invisible presence in the material world.

As Inner Light increases, it not only gives a Princess more magical ammunition, it also represents her magic bleeding into her mental and physical self. Characters with high Inner Light can have physical prowess beyond the limits of mere humanity. High Inner Light also makes a character more vulnerable to Sensitivity. 

	\bigsidebarstart
	
	\subsubsection*{Inner Light}
	
    \centering
    \begin{tabular}{|c|c|c|c|c|}
        \hline
        Inner Light & Ability Max & Wisps / Turn & Charm Max & Sensitivity \\ \hline
        1           & 5 / 7       & 10 / 1       & 3         & 1 die       \\ \hline
        2           & 5 / 7       & 11 / 2       & 4         & 1 die       \\ \hline
        3           & 5 / 7       & 12 / 3       & 5         & 2 dice      \\ \hline
        4           & 5 / 7       & 13 / 4       & 6         & 2 dice      \\ \hline
        5           & 5 / 7       & 15 / 5       & 7         & 3 dice      \\ \hline
        6           & 6 / 8       & 20 / 6       & 8         & 3 dice      \\ \hline
        7           & 7 / 9       & 25 / 7       & 9         & 4 dice      \\ \hline
        8           & 8 / 10       & 30 / 8       & 10        & 4 dice      \\ \hline
        9           & 9 / 10       & 50 / 10      & 11        & 5 dice      \\ \hline
        10          & 10 / 10     & 75 / 15     & 12        & 5 dice      \\
        \hline
    \end{tabular}
	\bigsidebarend

\noindent \textbf{Ability Max}: The highest a regular attribute or skill may be raised using XP. The number before the slash is the maximum level of any attribute or skill in normal form. The number after the slash is the maximum, including bonus dots, from Transformation.\\

\noindent \textbf{Wisps / Turn}: The maximum number of Wisps a character may have available in their pool, and the maximum number they may spend in a single turn.\\

\noindent \textbf{Charm Max}: The highest level at which a Charm may be invoked. Basic Charms do not go beyond 5 dots, but a Charm's effective rating is increased by one with every Upgrade. Inner Light caps the effective rating. The Princess may know more Upgrades to a Charm that would push the Charm above her Charm Max, but must choose which ones to apply when she uses the Charm. For Charms which apply a passive effect, such as manifesting Regalia, the Princess may change which upgrades are active with a Transformation action.\\

\noindent \textbf{Sensitivity}: The base number of dice rolled when a Princess's Sensitivity to suffering is triggered. \\

\subsection*{Wisps}
\phantomsection
\label{subsec:Wisps} 

A Wisp is perfectly mundane emotional energy. We often perform at our best in times of heightened emotion. For a Princess, her best is outright supernatural. Princesses use Wisps to fuel many of their abilities, and, like any human, she becomes tired if she runs low. A Princess's Inner Light determines her maximum amount of Wisps, as well as the maximum she can spend in one turn. 

\subsubsection*{Spending Wisps}
\begin{itemize}
\item    Charm Activation: While Transformed, a Princess may spend Wisps to activate Charms.
\item    Practical Magic: Whether or not she is Transformed, a Princess may spend Wisps to power Practical Magic, the abilities of which are determined by the Queen she follows.
\item    Holy Shield: While Transformed, a Princess may spend a Wisp to protect herself from one occurrence of damage from a external source, such as a single attack or one turn worth of damage from an Environmental Tilt. Aggravated damage becomes lethal damage, lethal damage becomes bashing damage, and all but the first point of Bashing damage is negated. Only one Wisp can be spent against any single source of damage. The attack could be turned aside from a vital spot at the last second, grazing her instead of striking home, or could be visibly stopped by a shower of sparks and light.
%\item    Open Heart: A Princes can speak directly to the heart, she places her emotions and feelings directly into her words. Anyone who hears them hears not just what she is saying, but what she truly means. It's not that useful for transferring facts or information, but it can be extremely effective at helping the other side see your point of view. If she attempts to to Force Doors she may spend a Wisp before rolling, on a failure roll Belief. Success on the Belief roll converts the Dramatic Failure back into a regular Failure, they may not agree with the Princess but they understand why she's acting this way and empathise with her. At least, they will if her emotions are relatable. If her motives are disagreeable or worse, the Storyteller should impose an appropriate penalty to the Belief roll.
%Editor: Skipping this commented-out part. 
\item    Quickened Transformation: A Princess may spend a Wisp to Transform reflexively without rolling. 
\end{itemize}

\noindent If the Princess's pool of Wisps ever falls to (10 - Belief) she gains the Condition Running on Fumes.

\subsubsection*{Regaining Wisps}
\begin{itemize}
\item    Call of Duty: A Princess normally regains Wisps by going out and making the world a better place in a way that fits with her Calling. At the end of every scene where the Princess successfully performs a task appropriate to her Calling, the Storyteller may call for a roll with the same Attribute and Skill as the task (if there is more than one applicable roll use your best judgment, supernatural bonuses and penalties do not apply) penalized by Shadows and grant one Wisp per success. If the Princess upholds her duty in a way that honors her Queen’s philosophy she gains one additional Wisp from the task, provided she scored at least one success.
\item    Inner Strength: When the chips are down and the situation is dire, a Princess can call upon her own inner reserves of determination. As an instant action, she spends a Willpower point and rolls Belief; if the roll succeeds, she regains Inner Light Wisps plus one Wisp per success. An Exceptional Success grants her next roll the Rote Action modifier. This ability can only be used when the Princess is in immediate danger, such as in combat, and is not meant to be an everyday means of regaining power - a good rule of thumb is, if it doesn't matter that this action is Instant and thus takes up a character's action for the turn, it's probably not dangerous enough to work. If the Princess makes a speech asserting herself and her Beliefs she gains 9again to the roll, with the Royal Tongue she may do this as part of the same instant action. 
\item    Circle of Light: Princesses with mundane families and friends often find that spending time with them, sharing their troubles or happiness, refreshes them for the supernatural fight. The Circle Merit represents these relationships; see that Merit for a full explanation. 
\item Voluntary Work: During downtime a player can say that her character spends her time fulfilling her Calling duties, in-between school, family or other commitments. State how your character fulfills her Calling and roll an appropriate Attribute + Skill - Shadows, each success grants one Wisp. If the Princess has the Mandate Merit a Successful roll grants one additional wisp per Mandate dot.
\end{itemize}

An Exceptional Success on any roll to regain Wisps removes a dot of Shadows.

\subsection*{Dreams (Aspirations)}
\phantomsection
\label{subsec:Dreams} 

All people aspire to great things, it's in our nature. For the Nobility it's not in their nature, it is their nature. Her Belief in making a better world is what gives her strength, it's what powers her magic, and it can consume her every waking moment. This desire to make a better world is represented by a Princess' Dreams. 

Dreams are a form of Aspirations unique to the Hopeful. In it's simplest form, a Dream is an Aspiration, where the primary purpose is to help somebody else, in a manner fitting the Princess' Calling. For a Seeker getting an A+ on the science test is an Aspiration, helping your friend get an A+ is a Dream. Winning a noble prize is an Aspiration, creating a cure for cancer (if it is motivated by saving lives) is a Mender's Dream. 

A Princess has a total of five Aspirations and Dreams, two of which must be Dreams and the remaining three can be either. Her Dreams provides all the same functions as an Aspiration, and provides one additional bonus. Whenever an Aspiration would provide a Beat a Dream will provide one Beat and one Luminous Beat. 

\subsection*{Luminous Experience}
\phantomsection
\label{subsec:LuminousExperience} 

Princesses work harder than most people. They have to if they want to improve the world. As a result, she also learns faster. This advantage is represented by Luminous Experience. 

Luminous Experience is earned in one of two ways. First, a Princess's Dreams provide Luminous Beats in addition to regular Beats; a Princess grows the fastest when she focuses on others. The second way is through volunteer work. When a Princess rolls an Exceptional Success to regain Wisps through volunteering, she gains one Luminous Beat.

In any situation where you can only gain a finite number of Beats, track Luminous and regular Beats separately. 

Five Luminous Beats form a Luminous Experience, which may be spent on Inner Light, Belief, or removing Shadows at the usual rate. (Mortals with Luminous Experience may spend it on Integrity.) Luminous Experience may also be spent on Social Merits, which represent the friendship and gratitude of the people the Princess helped earning that Luminous Experience. 

When calculating Experience costs, one Luminous Experience is equivalent to one regular Experience. Luminous and regular Experience may be combined for an expensive purchase. 

\subsection*{Belief (Integrity)}
\phantomsection
\label{subsec:Belief} 
\noindent Being a Princess is not something that happens at random. Princesses are here on Earth with a purpose: to make the world a better place and bring the Light of hope to everyone suffering. The existence of the Kingdoms proves that it can be done, and a Princess's conviction is one of her greatest assets. It is a source of strength: It keeps her going when times are bad, and it's what fuels her magic. If her Beliefs slip, as doubt creeps into her mind, her Inner Light within is beset by the darkness of the world.

Belief tracks the strength of a Princess's convictions. It replaces the Integrity system from \textit{God-Machine Chronicle}. Where Integrity measures how a character copes with traumatic and supernatural events, Belief only represents a Princess's ability to cope only with trauma and pain. (Since they themselves are supernatural, Princesses don't need to cope with the supernatural.) Belief is the core of inner strength and hope that allows a Princess to carry on even while she is suffering more than anyone could bear. and Belief works hand in glove with the Sensitivity system, which represents how much emotional pain a Princess is dealing with. 
%Editor: Revised definition of Belief in the second sentence and added an aside sentence. Does this still capture the intended meaning? 

Characters with high Belief are confident that, no matter how dark today may seem, things will be brighter tomorrow. When they face seemingly insurmountable problems, their principles shine like a lighthouse in their heart to show a path forward; they believe that if they remain true to their ideals and keep fighting, things will work out for the best. This makes Princesses with high Belief literally superhuman in their ability to cope; between Sensitivity and the horrors of the job, they have to be. A nine-year-old Princess might cry and hug her parents after seeing her brother's mangled body in a wrecked car, no amount of Belief prevents her from feeling, but that pain won't hold her back until she's pulled everyone from the wreckage, provided medical aid, and organized the bystanders.
%Editor: I wasn't sure I understood what the last sentence meant, and I left it untouched. (If nothing else, it is a run-on sentence.) I think it means "A Princess might feel very sad upon witnessing a personal tragedy, but she can contain the pain until after the emergency is over." Is that right? I can tweak the sentence if so. 

Characters with low Belief are wrecked with doubt and despair. They see the problems of the World of Darkness but no longer see a way to a better world. Her former ideals are now, in her eyes, naive, ineffective, or even harmful. At low Belief, a Princess's ability to cope is no more than the average person's, and that is not enough for a Noble who must bear the weight of the world on her shoulders. Without a clear path forward, her deeds are marked by desperation, and she begins to consider solutions she'd have once rejected as immoral.

As Belief falls and rises, players should consider how their character's views change. If a character \textemdash whose Beliefs include typical Western views on fair trials \textemdash magically acquires proof of a crime and administers vigilante justice because magic isn't admissible in court, her actions call her Belief in justice into question. 

This conflict is represented by a Compromise roll. A failure, and the lost dot of Belief, can represent her now feeling that fair trials are unnecessary or impractical until the world has improved. Future acts of vigilante justice will likely not trigger Compromises, though a material difference, such as a lower standard of evidence or a more brutal reprisal, might still be a Compromise. This Compromise might also inure the Princess to other Compromises about respecting the law; she might find interfering with a police investigation more palatable, for example. 

Even if the Princess succeeds in her Compromise roll and does not lose any Belief, she might still change her views. She might simply feel her actions were wrong and resolve to behave differently next time. She might instead adjust her Beliefs, creating a new but equally idealistic framework that contextualizes her actions and sets new checks and balances for the future.

\subsubsection*{Compromises}

Moral Compromises threaten a Noble's Belief. Belief comes in three tiers: paragon, idealist, and pragmatist. A Princess is affected by Compromises in her tier and below. If a pragmatic Princess considers an act unconscionable, she would reject it even more strenuously if she were an idealist. Each tier lists the dice pool for a Compromise roll, and it should be understood that this refers to the tier of the act, not the Princess. A paragon whose actions are a Compromise at the idealist level rolls 3 dice, not 5. As always, a player may choose to downgrade a failure to a dramatic failure in exchange for a Beat. 

Willpower cannot increase the dice pool for a Compromise roll, and magic cannot modify the roll unless explicitly stated otherwise, but some situational modifiers may apply. If the Noble Compromises in pursuit of her Vice, she subtracts one die, as her selfish behavior demonstrates disregard for her ideals. If she acts in pursuit of one of her Virtues, she adds a die, to represent her trying to navigate an ethical conundrum. 

If acting against her principles results in a positive outcome for the Princess, she loses a die, as this success encourages similar unprincipled behavior. If her actions cause trouble that principled behavior could have avoided, this reinforces the value of her principles and adds one die. A support network can help a Princess process dilemmas in a healthy way: A dot in Circle gives an extra die to Compromises at the idealist and paragon level. At four dots, Circle also provides an extra die to Compromises at the pragmatist level. Finally, if the Princess has a Condition caused by using the Invocation of Specchio, she gains one die; it becomes easier to protect her idealism through doublethink and denial. 

There is a another way for a Princess to lose Belief, and that is through draining her spirit with the magic of Lacrima. This works like a normal Compromise roll, but with a different dice pool. The details can be found in the section on Lacrima on page \pageref{LacrimaRef}. 

\begin{myindentpar}{10pt}
\emph{Dramatic Failure:} The Princess's faith in her ideals has been damaged. She is convinced that an ideal or principle she valued is unworkable or harmful and scorns her younger self for being naive enough to believe in it. Lose a dot of Belief and choose from the following Conditions (or create a new one with Storyteller approval) as her negative emotions send her magic haywire: Shattered Confidence, Dissociation, Scars of the Past, Abdicated, or Madness. Alternatively, you may take one of: Doubting, Heavy Crown, Self-Doubt, Obsession, or Over-Compensatory Guilt as a Persistent Condition. Either way, take a Beat. For a third option, and only when it would be a serious inconvenience, lose your powers for the remainder of the scene and take one of: Heavy Crown, Self-Doubt, or Over-compensatory Guilt, and take two Beats. 
%Editor: Would you be interested in putting the Conditions into bulleted lists? It might look weird, but it'd make this paragraph much more readable, especially at a glance. 

\emph{Failure:} The Princess's sense of hope is shaken. She is convinced that an ideal or principle she valued cannot work in this cruel world and is willing to put it aside until the world is ready. Lose a dot of Belief and choose from the following Conditions (or create a new one with Storyteller approval): Befriendable, Heavy Crown, Self-Doubt, or Over-Compensatory Guilt. Alternatively, and only when it would be a serious inconvenience, lose your powers for the remainder of the scene and take a Beat. 

\emph{Success:} The Princess comes through with her sense of idealism intact. Either she believes that she was wrong to go against her moral code and resolves to do better next time next time, or she adjusts her Beliefs to incorporate new life experiences, resulting in a different but equally idealistic moral code. However, she will still feel guilty or vulnerable during this period of transition. Choose one of the following Conditions (or create a new one with Storyteller approval): Befriendable, Heavy Crown, Self-Doubt, or Over-Compensatory Guilt.

\emph{Exceptional Success:} The Princess's experiences taught her the value of both her principles and pragmatism and how to best combine them in a complex world. She takes a Beat and gains a point of Willpower. \end{myindentpar}

\subsubsection*{Belief Tiers}

\noindent\textbf{Belief 8 - 10, High / Paragon (6 dice)} \\

\noindent A Princess at this level acts like she lives in the world of her dreams. A Princess of Diamonds might provide tuition at any time to anyone who asks, while a Princess of Spades might personally sneak into government offices to chat with officials without regard to social norms or trespassing laws. This is not to say these Princesses are naive fools. Maintaining this level of Belief for any length of time requires a Princess to be both skilled and astute. The Princess of Diamonds would need to be a fast and effective teacher with great time management, and the Princess of Spades would need to leave every official too impressed with her to press charges. 

Some paragons believe that the world is simply the product of our actions: it cannot improve until people behave like the citizens of a better world, and someone must make the first step. This view is popular in the Courts of Clubs and Swords. In the Court of Hearts, paragons warn that seemingly sensible Compromises can become embedded in societal systems and say prevention is the best cure. 

The World of Darkness is not kind to such idealism. Paragons must roll for Compromise when they don't live up to their own standards. But they must also roll when 
they need to question whether they can live their ideals in this World of Darkness; for example, a paragon could have failed an important task because her ideals held her back. Violating the third Oath of one's Calling is also a Compromise for paragons. Any successful paragon is sure to have faced several Compromises and refined her Beliefs until she has created a code of behavior that, with all her magic and inner strength, she can make work.

\noindent\textbf{Belief 4 - 7, Medium / Idealist. (4 dice)} \\

\noindent At this level, a Princess acknowledges the kind of world she lives in and has made concessions to practicality. She has given up on some of her Beliefs and now calls them unworkable or even a bad idea, but she retains her core ideals. For example, a Princess of Hearts who sees herself as a knight in shining armor may now feel ideals like courtly love and formal challenges to combat are impractical. 

However, she would still be a chivalrous knight: she fights only in tournaments or when necessary to defend the weak, and she inflicts minimum harm. A Princess of Spades might tell her fans to put on a school uniform because an education is worth a little conformity, but give different advice when wearing that uniform would require someone to suppress their sense of self. 

Compromises occur at this level when something challenges not only the practicality of a Princess's ideals, but their underlying principle as well. If violence broke out because the Princess insisted on integrated the two communities, this would be a compromise only if the Princess was a paragon. If a Princess of Clubs fails to prevent violence from breaking out because she insisted on integrating two communities, this would only be a Compromise if she is a paragon. But if she came across two groups living peacefully by ignoring each other and her attempt to integrate them caused violence, or the integration results in something less than either of the original two groups, then the Noble is forced to ask if integrating different groups is a good idea in the first place. Such a Compromise would fall under this level.
%Editor: I don't understand the difference between the paragon example ("If a Princess of Clubs fails to prevent violence...") and the first idealist example ("But if she came across two groups..."). Both seem to have a Princess's effort to integrate two communities cause violence. Please clarify. I'll make only small edits to this paragraph for now.  

If a Princess demonstrates disregard for her ideals, then a Compromise roll is appropriate. A Princess could, for example, give up on integrating two groups for exam prep. Breaking the second oath of one's Calling also lands here. 

Inflicting severe emotional pain, manslaughter, and other crimes of a similar level usually fall here, with obvious caveats. Sending a murderer to jail might inflict emotional pain upon him, but that would trouble few Princesses' Beliefs. 

Each Princess has different ideals. A Princess of Diamonds's strategy might cause someone's death, and she would have to ask herself if getting everyone to think things through like she did is enough to make a better world. Or if a Princess of Spades's prank leaves the pompous but well-meaning official in tears; and is forced to reconsider if being an authority figure who makes bad decisions means it's automatically OK to take you down a peg or two. If a Princess inflicts real harm to another person, her Sensitivity needs her to ask how she made such a mistake, and the Princess will question if her underlying principles were flawed. A Champion who kills her opponent in battle is no exception. Magic provides so many non-lethal options that a modern day samurai who fails to take a human opponent alive will have to ask herself if Bushido can make her into someone who can wield violence without inflicting undue harm. However, killing creatures of the Darkness (apart from Darkened) is never an issue, as such monsters are avatars of evil. They have no moral worth and cannot be redeemed. 

%Editor: I don't understand this paragraph. The examples of the Diamonds and Spades Princesses don't seem to match with the first sentence. The first sentence is about differing ideals, while the examples seem to be about Princesses causing unintended harm. The last example, about the Champion, is phrased in such a way that it seems to be comparing itself to the other examples. I guess I'm asking for clarification on what the main point of this paragraph is, so we can focus it down. For now, doing only some light edits to it. 


\noindent\textbf{Belief 1 - 3, Low / Pragmatist. (2 dice)}\\

\noindent When a Princess has been beaten down, she might decide to just let the world win. She tried idealism. It didn't work. Now, she'll do whatever's necessary. This is not to say she is happy about it: A pragmatist Princess's actions pain her greatly. Nor does it mean that she has no principles beyond mere pragmatism. Though she has given moral ground, the Princess will cling very tightly to the principles she has left. A Princess of Storms will cut through innocents to destroy Darkspawn, but she would never, under any circumstances, aid a servant of the Darkness.

Compromises at this level require more than a challenge to a Princess's ideals. If a Princess has sunk to this level, then she would have faced so many challenges that she would be inured to further questions. To face a Compromise at this level, a Princess has to actively act in opposition to her own principles. No matter how many times she's justified her actions, opposing her own principles creates a challenge she cannot ignore. Similarly, abandoning a principled goal for selfish reasons like money falls at this level. Most pragmatists and all idealists will face a Compromise roll for torture, murder, and other heinous crimes, as such actions are in opposition to most principled behavior. Finally, breaking the first oath of one's calling lands here. 
%Editor: Is "Oath" something you folks want to capitalize? Also, in the second-to-last sentence, "Most pragmatists and all idealists" implies that paragons DO NOT need to make a Compromise roll for stuff like torture and murder. I'm assuming that's not intended? 

\subsubsection*{Worked Example}

Polly P. Peterson is a Seeker of Hearts who works as a journalist. As a Princess of Hearts, she believes strongly that authority belongs to those who earn it by winning the support of the people through virtue and competence. She also believes that this system cannot work without good journalism. Without accurate and unbiased news, the public cannot make informed decisions about who should run their society. Without a public mandate, there can be no governance. Without governance, society falls apart. For simplicity, this example will only look at Polly's Beliefs that relate to journalism, and it will treat Polly's Beliefs as static, instead of things that will grow and adapt during play. 

Paragon
As a paragon, Polly would roll for Compromise every time she failed as a journalist. When she notices her bias leaked into an article that's already gone to print, or when she learns a slick politician's lies got past her into the morning edition, the dice must roll. When Polly's experienced but cynical mentor says she needs to spend more time on money-making celebrity gossip and less on informative political news, that alone would not be a Compromise roll, but if she took his advice or if Polly's insistence on focusing on reporting real news caused her a significant setback such as failing to get the job at a prestigious paper she wanted, that would be a Compromise. To succeed as a Paragon, Polly would have to be the ideal journalist. She would need a level of professionalism and competence that only Nobility could achieve. 

Idealist
At the idealist level, Polly would roll when something challenged the idea that good journalism leads to an informed public making good decisions. Polly would no longer risk her Belief by making a mistake or writing a truthful gossip piece to pay the rent. But if she saw that none of the major papers were reporting on an important issue, Polly would see that as a worrying disregard for a journalist's core purpose. If Polly's report on an important issue made the problem worse — say, by spurring an angry public into the hands of a demagogue — that would need a Compromise roll. 

A long career as an idealistic journalist is within reach of any Princess. Most circumstances that could trigger Compromises are the results of Polly's own choices and within her power to avoid, and the exceptions are rare enough that good dice rolls or the occasional Experience expenditure can keep Polly's Belief stable, though collaborating with mortal institutions or other Nobles to ensure demagogues aren't the loudest voices responding to her expos\'e cannot hurt.  

Pragmatist
As the pragmatist level, Polly would make a Compromise roll if she subverted the public's ability to get accurate news or journalism's ability to provide that news. She can abuse her position as a journalist to fulfill her royal duties, but this would \textit{not} require a Compromise roll. For example, she could pretend her investigation found nothing, if the corrupt executives agreed to demolish a Tainted building and put affordable housing in its place. Or she might investigate a political candidate she dislikes and publish every skeleton she finds in her closet while writing nothing about the candidate she supports. Sabotaging another journalist's investigation would require a roll, however, as would intentionally spreading falsehoods to the public via journalism. Polly, like most Princesses who fall to the pragmatist level, can only fall further through her own actions.  
%Editor: I think the three paragraphs in this section are long, unapproachable blocks of text. The section could do with some further dividing. Is there such a thing as a subsubsubsection? Was thinking we could put a heading on each one of the levels, and then we break each level's paragraph into multiple paragraphs (so that each level has two or more paragraphs). That would make the section an easier read. I've included the headers I had in mind, and I'll my best to format them. 

\begin{mybox}


\subsubsection*{Mean Girls}

A Princess's magic flows from her Belief. Destroying her self-confidence can be an effective way to weaken her powers and a devastating precursor to a physical attack. The rules for Social Maneuvering (\textit{Chronicle of Darkness}, page 81) can be applied for this purpose.

The aggressor sets the goal of \textquotedblleft inflict a Belief Compromise" or \textquotedblleft trigger a Sensitivity check." Calculate the Doors as normal. If the aggressor is targeting Belief, this always counts as opposing a Virtue. When the last Door is opened, make a Compromise rol or Sensitivity check as normal. The Princess may not offer an alternative. Vices and Leverage can work as normal; tickets to a Princess's favorite band could make her spend an evening with her tormentor, who can use that time to challenge the Princess's ideals and chip away at her self-confidence. 

Targeting a PC requires cooperation from a player. If someone is intentionally trying to ruin a Princess's Beliefs and self-esteem, that's a very good justification for the player to say her character's impression just reached hostile. Remember though, there are Beats for going with the flow ((\textit{Chronicle of Darkness}, page 82) and Compromises.

Princesses can spar verbally and physically at the same time. When they do, both parties make a contested social roll every turn of combat, with the winner opening a Door. Once again, player consent is required. Princesses are defined by their strong Beliefs, so it would be very unusual if they were constantly at risk of throwing in a fight because of a few well-placed words. But when the circumstances are right, it is very fitting for the genre for a magical girl to defeat (and perhaps befriend) her rival with words.

Triggering Sensitivity doesn't always require Social Maneuvering. Taking advantage of an appropriate Condition, like Embarrassing Secret, or doing something emotionally shocking, like executing a hostage, might be enough to trigger Sensitivity there and then. This depends on the severity of the Condition and the emotional resilience of the Princess.
%Editor: I changed the capitalization of some terms in this sidebar to match how they look in Chronicles of Darkness. 

\end{mybox}

\subsubsection*{Bad Example}

Most moral philosophies, and most Princesses, agree that acting immorally through a proxy is still acting immorally. If it is wrong to murder, then it is wrong to ask someone to murder on your behalf, and asking would be a Belief Compromise, as though the Princess held the knife herself. 

But even if the Princess doesn't ask directly, her Belief may still be in danger. If her follower violates her principles while acting on her behalf and could reasonably think that the Princess would approve, then the Princess is forced to consider whether her behavior has fallen short without her realizing. 

\textit{For example: Princess Rosa, a Champion of Hearts with Belief Six, believes she must always attempt to resolve a conflict peacefully before using violence. Somewhere down the line, she concluded that it is acceptable to use her Emotional Intuition to check how a character will react to an attempt to resolve things peacefully. If they will not respond reasonably, she can save her breath. (This would be a Compromise, but in this example, Rosa succeeded on that Compromise roll and modified, rather than abandoned, her code). However, Rosa's followers, a neighborhood watch group she founded, do not have the Emotional Intuition Merit, so when Rosa discovers that they've been following her example, only with mundane Empathy rather than supernatural intuition, she realizes that her code of ethics is unstable for her position and must make a Compromise roll.}

Storytellers should apply this example sparingly. It should only lead to a Compromise roll when the Princess's actions leads to bad behavior that would not have otherwise happened \textit{and} when the link between the Princess's behavior and her subordinates' behavior is clear. 


\subsubsection*{Effects of Belief}

% The state of a Princess' soul shines outwards onto those she holds in affection, an ability refered to as her Echo. Hopeful with a high Belief inspire confidence, and sharpen the voice of conscience; at Belief 8 or more, a Princess grants mortals (anyone using the default Integrity trait) who are in regular contact with her a +1 bonus to all Resolve-based rolls when they are trying to remain true to their principles or morals, this includes all Breaking Points except ones caused by the mortal breaking their own morals. Contrariwise, Hopeful with a dwindling Belief sow doubts and deaden consciences among her acquaintances; a Princess with a Belief of 3 or less inflicts a -1 penalty to all Resolve-based rolls made by mortals in regular or prolonged contact with her. It takes several hours of peaceful interaction, for good or bad, before a mortal begins to be affected by a Princess' Echo. 
%Editor: Was going to ignore this commented-out section, but I see that it gives useful info for Effects of Belief. The definition of Echoes is the biggest thing; it hasn't been defined since the introduction. Not sure what you wanted to do with this commented-out part, but, if nothing else, integrating parts of this paragraph into the rest of the section would be helpful. I'll leave this part untouched for now and won't edit out mentions of Echoes in the rest of the section, but this should be revisited. 

As Belief rises, the Hopeful's inherent magic glows more brightly and becomes more potent. Her inner, emotional strength becomes magical strength. In a conflict of opposed magic (see Clash of Wills, page 110), a Princess's Belief becomes part of her dice pool. A high Belief Noble radiates confidence and optimism, increasing the odds she'll inspire people through a positive Echo. High Belief also enhances a Princess's physical strength. 
%Editor: I suggest linking to the section on Clashes of Wills on page 110. Tried to make the link myself, but I might have gotten it wrong. Lemme know your thoughts about this idea. 

Characters with low Belief find themselves lacking the strength to bear the world's pain. As Belief falls, the lash of Sensitivity becomes ever more threatening, and her doubt begins to inspire similar feelings in those around her, increasing the risks of spreading a negative Echo. 

\subsubsection*{The Enlightened at Belief 0}

\begin{quote}I used to think that you were a guardian angel, come to answer our prayers… But Lucifer was an angel too, wasn’t he?\\
\textemdash\, Professor Emil Hamilton, \\ \textit{Justice League Unlimited} \end{quote}

When she Blossoms, the Princess feels a strength of optimism never known to those who have not touched the Light. Nothing seems impossible. With hard work, dedication, and conviction, every problem can be solved. No one needs to suffer anymore. Only the most sheltered Princess can keep such attitudes for long: Belief and optimism must give way to reality. The Princess will learn that real problems can't be solved in an afternoon, that scars take more than a impassioned speech to heal, and that she herself is limited. She does not have the time or energy to fix every single problem.

Most Princesses remain here for most of their lives. They have compromised their impossible idealism in the name of practicality, yet retain a strength of conviction few can match. Some fall much further. If darkness is the absence of light, the absence of a Princess's Light is very dark indeed.

When a Princess hits Belief 0, the anchor of her self is shattered. As her sense of hope, drained by so many Compromises, runs out, her Inner Light finally dims. She decreases her Inner Light to 0 dots — nothing more than a few fading embers. Her Phylactery, the item Princesses need to Transform, rapidly rots away from age, and the Princess is forcefully de-Transformed. She is barred from the Dreamlands; the Queens can offer no comfort. Even the Twilight Queens are disgusted by a Noble who falls this far: Tears sees her as a threat to Alhambra, Storms as just another Dark creature waiting to happen, and Mirrors as a failed True Queen.
%Editor: Phylacteries haven't been defined yet in the text (excluding the intro). I revised the paragraph slightly to define Phylacteries, but I'll refrain from linking to their definition, as it seems distracting in this context. 

Many would say that this is enough of a curse, for, stripped of rank and privilege, she is left to wander the Earth. But that is not all. She would not not merely be a fallen Princess, but a fallen human: little more than a shuffling husk that drifts through life, barely recognizable as a former champion of the Light. Stripped of all drive, she might just lie still until she starves, unable to see the point of eating. If she is pushed, she would go about her life, but she would be too empty inside to say no, to say that she hurts too much. She would be a hollow shell of a woman, lacking the strength to fight for what she believes in, always folding at the least pressure, unable to even cry for help. 

When a Princess drops to Belief 0, the final act that broke her Beliefs inflicts the Soulless (\textit{Chronicle of Darkness}, page 290) and Enervated Conditions. For as long as her Belief remains zero, she cannot access any of her powers from the Princess template. Since the Princess's soul is not actually missing, just falling apart, finding a replacement soul cannot heal her. The only solution is for the Princess, with the help of her friends, to rediscover her Belief before she loses her last dot of Willpower. Her friends can do this using the Social Maneuvering rules [\textit{Chronicle of Darkness}, pages 81-84]; the number of Doors they must open equals twice her Shadows dots. If all the Doors are opened, and the Princess spends 2 Experiences and sacrifices a Willpower dot, she can remake her Phylactery, return her Belief to 1, and end the Soulless and Enervated Conditions.

A soul-dead Princess who does not get help from others faces a final choice when her last dot of Willpower fades: to die as a mundane human, or to yield to the Darkness and rise as a horror. To determine which, the Storyteller rolls the Noble’s Shadows as a dice pool. 

If the roll fails, the Princess drops into a coma from which she never wakes - her heartbeat ceases within a day at most. 

If it gains even one success, the former Princess attempts to Transform despite her lack of a Phylactery. She turns toward the Dark to fill the empty void where she once held her Light. She focuses on her grief and bitterness, for in her deranged state, even that is better than nothingness. The mad will of a Princess in this state is a terrible thing to behold. She laughs with pure grief, she cries tears of crystal clear despair. Her feelings echo around her, calling pure Taint up from the Dark World and shaping it into twisted reflections of her own life. She even calls the Taint into her own body, fueling her final Transformation into a being drawn from the depths of her own nightmares.
%Editor: Splitting this paragraph up to make it easier to read. Also, the paragraph contradicts itself about whether embracing the Darkness is a choice. The first sentence says that it's a choice, but the mechanical explanation makes it sound like it's down to just a dice roll, with the player not having any input. Please clarify. 

Transforming like this should not be possible. But those who do it anyways are no longer Enlightened or even human. They shed their mortal forms, as Darkness rushes into the void the Light once filled, and they are both twisted and empowered by this blasphemous process. They wear blackened crowns on monstrous heads; they become the Dethroned.

\subsection*{Sensitivity}
\phantomsection
\label{subsec:Sensitivity}

Being born from a shining, empathetic, heart, a Princess is by her very nature‭ emotionally ‬vulnerable. Emotions strong enough to destroy Darkspawn and outrun cars flow through her. When those emotions turn against her, the consequences are debilitating‭ at best. In the wrong circumstances, they can be deadly. When cultists kill a hostage in the middle of battle, the horror paralyzes her and leaves her open to attack. When her sadistic teacher bullies a classmate, every flinch feels like ice in her veins, and she is compelled to act, even though she knows she can do more good in private after class. All magic comes with a price, and this is a Princess's. She will pay it all her life. Neither age nor experience will free her, for her emotions only get stronger as her Inner Light grows. Any attempt to develop emotional armor is futile. Her own magic breaks such walls open to let her emotions get out and affect the world. 

But the Nobility don't spend their entire lives one bad breakup away from spontaneous combustion. A Princess's positive emotions are just as superhumanly potent, allowing her to shoulder grief that would crush anyone else. Her eternal hope and optimism is her bulwark against doubt and heartbreak, the foundation of her emotional and magical stability. 

Whenever the Princess first enters a state of significant emotional pain or turmoil, she makes a Sensitivity roll, with the dice pool defined by Inner Light. Deliberate acts of cruelty are especially shocking to the Princess and add an extra die to the roll. This increases to two dice if she is intimate with the villain or victim, or three dice if she performed the action herself. The Storyteller may add up to three additional dice if the act is especially heinous or emotionally shocking. 

When a Noble's attention is drawn to someone else's emotional pain, the effects can be worse than her own problems. Her Sensitivity means she feels she cannot help but feel their pain herself. More than that. Sensitivity rolls made because of a mortal's emotional pain gain the 9 Again modifier. If that someone is an innocent (such as a child), the modifier rises to 8 Again. If a single action inflicts emotional pain upon multiple people, the stronger modifier applies.
%Editor: The second "Sensitivity" in this paragraph used to say "Sympathy," as in "Sympathy rolls made because of a mortal's..." From the context, it seemed like it meant "A Sensitivty roll, made in sympathy to a mortal's pain." Calling it a sympathy roll seemed confusing, so I changed it. If a sympathy roll is a specific thing, please change it back and let me know. 

Even the Enlightened's empathy has limits. If a Princess witness a sadist in great emotional pain because she rescued his victim, her empathetic nature means she would sense that pain. However, the pain would be so alien to her that it would be easily compartmentalized and ignored. Many Princesses compare this to reading about someone's feelings in a book or feeling butterflies ten feet outside their stomach. When this happens, no Sensitivity roll is required. The Storyteller is the final arbitrator of whose suffering can trigger Sensitivity. 

\begin{myindentpar}{10pt}    

\emph{Dramatic Failure or Failure:} The Princess retains control over her emotions.

\emph{Success:} The Princess is staggered. She loses her Defense and cannot act for the next turn. If she is Transformed, she loses a Wisp in a random magical discharge. In addition to any other effects, the Princess gains both a Dot of Shadows and the Fresh Shadows Condition. If the Storyteller and player agree, she may also take the Girl Underground Condition. 

\emph{Successes/2 (rounding down) + Shadows is equal or greater than Belief:} Take one of the following Conditions (overtly magical effects require her to be Transformed) or, with the Storyteller's permission, the player may choose or create a new one of similar severity: Abject Denial, Primal Scream, Shadow Takeover. The Condition should represent the Princess losing control over her actions or her magic or the consequences of said loss of control. If an emotional shouting match will have serious consequences, then a Condition representing those consequences is a valid choice. Alternatively, when it would be a serious inconvenience, the Princess may take a Beat and lose all her magic for the scene, reverting to her mundane form. In addition to any other effects, the Princess gains both a Dot of Shadows and the Fresh Shadows Condition. If the Storyteller and player agree, she may also take the Girl Underground Condition. 

\end{myindentpar}

\begin{mybox}



\subsubsection*{Sensitivity and Compromises}

The events that trigger Sensitivity and provoke Compromises are often one and the same. A Princess who believes in both the power of romance and healing the sick could face both Sensitivity and a Compromise when her boyfriend leaves her because she's been spending too much time at the hospital. To uphold one principle, she must compromise on the other, and a sudden breakup is an emotional shock.

However, that same Princess would face only a Compromise if she came to that realization herself and she and her boyfriend agreed to end their relationship. A Princess may not be able to build emotional armor against the shock of a breakup, but an experienced Princess can see it coming and avoid the shock. 

Meanwhile, an entirely different Princess who didn't believe the power of love was integral to building a brighter future would only face Sensitivity after a sudden breakup.

If the same event triggers both a Compromise and a Sensitivity roll, roll for Compromise first.

\end{mybox}

\subsection*{Shadows} 
\phantomsection
\label{subsec:Shadows} 

Shadows are the painful product of Sensitivity. They are the if-onlys, should-have-dones, and could-have-beens that haunt a Princess. They're a cancerous train of thought that can spread its tendrils through her mind and force its way to the front of her thoughts at inappropriate moments. For Princesses, whose emotions are magic, they can be as debilitating as any illness. A Princess with Shadows upon her soul not only feels sympathy for the suffering she witnesses, she internalizes them. She suffers flashbacks of the moment that created her Shadows. If she couldn’t help then, why should she expect to do better now? These thoughts become self-fulfilling: Because she feels lesser, her magic becomes lesser. It can too easily become the beginning of a vicious cycle: the Princess believes she cannot help, therefore her magic weakens, therefore she has trouble helping people, therefore she believes she cannot help.

It is fortunate then, that so many Princesses cultivate friendships and allies who can help them break the cycle.

Dots in Shadows impair a Princess's access to her powers in two ways. First, they are subtracted from her Transformation roll, which becomes Belief + Inner Light - Shadows. Secondly, Shadows impair a Princess's ability to recover Wisps: Her Wisp income from the Circle Merit is reduced by Shadows, and all the dice pools she rolls to regain Wisps subtract 1 die for each dot in Shadows she possesses.

When a Princess gains Shadows, she also takes the Fresh Shadows condition. If she finds closure before the end of the scene, she may remove the dot. After the scene ends, the dot \textquotedblleft sticks\textquotedblright\, and can't be removed so easily. If a Princess with Shadows fulfills her Virtue, she may remove 1 dot of Shadows in addition to regaining Willpower. If a Princess rolls an Exceptional Success on any roll to regain Wisps, she may also remove a dot of Shadows. Finally, a Princess may spend Beats to remove a dot of Shadows at a rate of one Beat to remove one dot of Shadows.  

%\subsection*{Using Sensitivity}

%Sensitivity does have it's advantages, or rather it has one small advantage. A Princess can use Sensitivity to read another person's emotions by rolling Sensitivity vs Composure + Supernatural Advantage. On a Success the Princess feels the targets emotions as though she were experienceing them for herself. 

%It's not a particularly effective system when compared to certian Charms or even mundane Empathy. However feeling another's emotions, rather than simply understanding them, makes it much easier to learn Charms from another Princess. It is also used to reinforce group bonds within Nakamas.
%Editor: Ignoring this commented-out part. 

\subsection*{New Trait: Transformed Abilities and Skills}
\phantomsection
\label{subsec:TransformedAbilitiesAndSkills} 
Next to each Attribute and Skill on a Princess's character sheet, in addition to the five blank dots normally present, there are another set of dots in parentheses. These dots represent additional Skill and Attribute dots the Princess has while Transformed. While Transformed, your Transformed Attributes and Skills are automatically added to your normal Attributes and Skills; this can change advantages such as Speed, Willpower, Initiative, and Health, so record these values for both your mundane and Transformed forms. 

If you take damage in Health boxes derived from Transformed Stamina, you may spend Wisps to heal that damage when you return to mundane form, per the rules for Holy Shield. If you don't do so, the damage wraps around, as usual for temporary Health. Willpower points derived from transformed Resolve and Composure can be spent only while you are Transformed. However, if you have spent them, you can regain them at any time. If you have spent both mundane and transformed Willpower, you regain the mundane Willpower first. If you sacrifice a Willpower dot, you lose a transformed Willpower dot first.

Your total Transformed ability may not exceed the maximum rating for Transformed abilities as determined by your Inner Light, and you cannot have more dots in Transformed abilities than you have Belief + Inner Light; Transformed Skills only count as one third of a dot when counting against this limit. A Princess who loses Belief loses Transformed Dots and is refunded the experience spent.

\begin{mybox}

\subsubsection*{Society and Skills}

A Princess practicing her mundane body will earn mundane Attributes and Skills. Training in her Transformed identity, in the Dreamlands, or by remembering abilities from past lives will create Transformed Attributes and Skills. Many Princesses seeking mundane skills prefer to train twice. Skills take the same amount of time to train, whether in a mundane or Transformed state. 
%Editor: I'm not sure I understood what the last sentence means and tried splitting it in two and rewording it. Is this the intended meaning? 

In terms of gameplay, this doesn’t have much effect. If you can only justify earning a Transformed dot now but want a mundane dot, you can usually fit one Experiences worth of training into the next downtime. 
%Editor: Previously, this said "experience point" instead of "Experiences." I changed it because it seemed like inconsistent terminology. Is "Experiences" the intended word? 

It's only when the Princess is constrained by external factors that she might be unable to transfer skills from her Transformed identity to her mundane one. Say a Princess wanted to train to be a better shot. In many parts of the world, even getting access to a gun is difficult for a grown adult, let alone a child. In such cases, the players will have to dedicate screen time and magic to overcoming these problems.
%Editor: Added sentence here to clarify what I assumed was the intended meaning for bringing up gun access. 

\end{mybox}

\end{multicols}\newpage\subchapquote{Champion}{\textit{When I started, I thought this would be some grand adventure, sneaking out with my best friends to fight monsters. I never thought any of us would actually die. But even after everything that's happened, I'd still choose this life, because somebody has to. Either it's me, or this burden is forced onto somebody else.}}
\phantomsection
\label{scha:Champion}\begin{multicols}{2}\raggedcolumns

\noindent Champions epitomize the ancient concept of a hero -- those who accomplish. In the Kingdom, they were the defenders and the conquerors, the dragon slayers and the chain smashers. Their counterparts in the Reborn world value strength for altruistic reasons. They serve the weak. Moreover, they enjoy physical accomplishment for its own sake -- Champions frequently sprout from athletes, those with a penchant for construction or day labor, and even brawlers. Anyone who loves to push their body to the limit or appreciates the simple need for a sturdy pair of hands makes for a grand Champion.
%Editor: "Reborn" shows up here as a proper noun, but it hasn't been mentioned yet in the text — not even the lexicon. I'm not 100% sure of the meaning, so is "mundane" an acceptable substitute? 

Given their emphasis on physique, some tend to mistake Champions as simple-minded or naturally violent. Rather, they just look for direct solutions. To them, problems are simply tasks that need doing. If they can't find a physical solution, something they can just \textit{do}, then Champions will break the situation down until they can. They cut the Gordian knot, instead of puzzling over it. If Champions have a weakness, it's that they tend to be blind to more subtle pain and motivations, but it also makes them refreshingly straightforward and realistic. 

\section*{Dreams of the Kingdom}

Champions dream of victories on a grand level -- armies defeated single-handedly, mountains knocked down, rivers wrestled from their paths, buildings held up for days at end. Their exploits are the stuff of Gilgamesh, Hercules, and Samson. But when they dream of the Fall, they see weapons snatched from their hands, their own bodies turning sickly and weak, and dark throngs overpowering and humiliating them. 
%Editor: Two questions a bout this part. First, the fall of the kingdoms isn't mentioned before this as a proper noun and rendered as "the Fall." Could we change it to something like "the fall of the Kingdoms"? Second, what is this paragraph (and future similar sections) describing? When Champions fall asleep, do they all dream dreams like the ones described here? If so, I think it'd be good to say that explicitly. One or two sentences at the start would do the trick. Thoughts? 

check about the Fall. also check Experience(s)? 

\section*{Magic}
Champions have affinity for the Bless, Fight, and Perfect Charm families.

Starting Champions get one Transformed Attribute dot in either Strength or Resolve. 

\section*{Duties}
Champions regain magic whenever they perform a task for somebody who is physically incapable or doesn't have the time and/or resources they need. This can mean doing chores for an elderly neighbor, building homes for the poor, protecting a classmate from a bully, or otherwise making a sacrifice of your time and effort for someone else's benefit. 

Champions can also restore their magic by protecting the innocent. Wiping out a nest of Darkspawn, defending a child from a wild animal, even walking the beat as a police officer: all of these can all fulfill a Champion's duties. 

\section*{Oaths}
\noindent \textbf{First Oath:} I am a protector, not a tyrant. I am given a sword to defend the weak and the innocent. Inflicting violence for personal gain betrays all that the Light has invested within me. \\

\noindent \textbf{Second Oath:} I must never abandon or ignore those who need my protection. If someone is in danger and I have the power to protect them, my power gives me the responsibility to step up to the challenge.\\

\noindent \textbf{Third Oath:} I shall swear myself to a code. Might is never the ends, only the means. My code shall determine when violence is acceptable, and for what reasons a sword may be drawn. I may change my code, but I must always have one. Once sworn, I must uphold the values and principles of my code. 

\section*{Stereotypes}
\begin{itemize}
\item Graces: By the time you have raised an army, I will have already won the battle. 
\item Menders: You tend to the wounded. There will be plenty of them soon enough.
\item Seekers: You want the truth? Ask a Seeker. You might not like what you hear, but you probably need to hear it. 
\item Troubadours: Why are you just standing there singing? Hit something! 

\item Vampires: We have a word for these where I come from -- \textquotedblleft targets."
\item Werewolves: Nice doggy... we're on the same side, remember?
\item Mages: Don't make me come over there.
\item Prometheans: I feel sorry for you, really, but you have to move on now.
\item Changelings: Yes, your life sucks. Now what are you going to do about it?
\item Sin-Eaters: If you narrowly escape death, but then make yourself a slave to a dead thing, you haven't really escaped death at all, have you?
\item Mummies: You can't remember why you're fighting, I wouldn't want to be you. Not for all your power and immortality.
\item Beasts: Spare me the bullshit. Are you going to stop hurting people or are we going to fight?
%\item Deviants: Stop running, stand behind me and we'll hold our ground right here. 
%Editor: Is the commented-out item something you want to restore? 

\item Mad Scientists: So, could you show me how to use that thing?
\item Leviathans: RUN! I'll buy you time... and goodbye. 

\item Hunters: You guys are crazy! How can I help?
\item Mortals: My job would be a lot easier if these guys would step up, but some of them just can't. That's why I'm here. 
\end{itemize}

%Editor: Would you consider changing the name of this section (and other similar sections) to something like "How Champions See Others"? 
\section*{Concepts}

Knight in shining armor, firefighter, beat cop, volunteer construction worker, inspiring general, wilderness guide, old-school warrior prince, swashbuckling adventurer. 

\section*{Inspiration}
Shoutan Himei, the Sailor Senshi, Nagisa Misumi, Buffy Summers, and the mag1ic warrior archetype in general. 

\end{multicols}\newpage\subchapquote{Grace}{\textit{Leadership is so much more than making the hard choices. To be a leader, you must represent the best in the people you serve, so that when you make those hard choices, they can see themselves in you the and find the strength to make that choice too.}}
\phantomsection
\label{scha:Grace}\begin{multicols}{2}\raggedcolumns

\noindent If we lost anything in the Fall, it was each other. People now wrap themselves in an shell of cynicism and distrust to save themselves from humiliation and risk. They carry popular opinions like armor, with no conviction. People can't save each other if they can't understand each other, and they can't save themselves if they don't even know who they are.

Graces are the fair maidens portrayed in song and art as a symbol of order in stability. In ages past, they were teachers, leaders, lawmakers, and judges who helped a species of almost-monkeys become something more. Now, though their duties have lessened, Graces continue them in the modern world. Graces force those around them to look at the people around them and confront what they find there. They are not malicious and they are not cruel – they know there can be pain, but they only want to show everybody what they can already see: the golden light that could shine in all of humanity.
%Editor: I'm not understanding the fourth sentence ("Graces force those around them to..."). Could you please clarify? Leaving that sentence untouched for now. 
 
Graces are people people. They like to know people, they like to hear about people, they like to be around people, and they like when people are people. Nothing pains them so much as someone denying their own humanity. Others would describe them as “intrusive” or “presumptuous” — always nosing into other people's business and offering advice when none is solicited. A Grace might be written off as naive or idealistic, for all her proselytizing about friendship. They do tend to be outspoken and emotional. They never restrain themselves from speaking what's on their mind or in their heart – and might be no strangers to embarrassment because of it. 

\section*{Dreams of the Kingdom}
Graces dream of fantastic journeys across the world, witnessing moments of beauty and love unique to civilizations they've never even dreamed of. On their travels, they celebrate all the uniqueness and humanity the world had to offer. They partake in enormous ceremonies to forgotten gods, and debate in enormous parliaments of wise and dedicated leaders. But as the wind turns cold and the stars fall from the sky, the people turn into silent shadows who drift about in a monochrome haze, and everywhere she travels is the same unending silence. 

\section*{Magic}

Graces have affinity for the Bless, Connect, and Govern Charm families.

Starting Graces get one Transformed Attribute dot in either Presence or Composure. 

\section*{Duties}
Graces recover their magic whenever they convince people to act on their virtuous impulses or bring them closer to their fellows. This covers situations such as convincing an addict to seek out help, easing the anger between antagonistic relatives, getting a classmate to study for her own sake instead of cheating, or even just helping someone act on their crush. 

Graces also refresh their magic by leading societies. This can include turning a petition into a law, making executive decisions in a company, debating in a student council, working as a judge, and making a convincing speech. All are sacred duties to Graces. 

\section*{Oaths}
\noindent \textbf{First Oath:} I must never force another to do as I wish. I can persuade and inspire but never blackmail, or worse, mind control, as that would betray everything I stand for. Nothing good can come from such rotten foundations. \\ %These should be the other way around, but this allows Storms to encourage the Vice of hatred. 

\noindent \textbf{Second Oath:} My purpose is to help people be the best version of themselves. I should always seek to uplift, never pushing anyone to give into their vices and worse impulses. \\

\noindent \textbf{Third Oath:} I must be honest, not manipulative. Mine is the magic to build trust and bring people together, not the power to turn people into tools. 

\section*{Stereotypes}
\begin{itemize}
\item Champions: What will you fight for? Me, perhaps?
\item Menders: You can only do so much by treating the symptoms. 
\item Seekers: Look as much as you want. I have no skeletons in my closet. 
\item Troubadours: I'll lead them, and you shall inspire them.

\item Vampires: Explain to me again how your civilization is anything of the sort. 
\item Werewolves: I might never know loyalty like you, but your pack has five. In a week, I could make fifty friends. 
\item Mages: Your magic isolates you. Nothing good can come of it. 
\item Prometheans: You want me to teach you how to be human? Well, you've come to the right person. 
\item Changelings: If you're not careful, a faerie can turn your words on you.
\item Sin-Eaters: This mythology was stolen from \textit{Dragonlance,} and it sounds like you wrote this code of honor while drunk! Is there nothing of substance to your society?
\item Mummies: A leader and a servant? I like the dichotomy. If only it wasn't so dysfunctional. 
\item Beasts: I've dedicated my life to helping people unlearn the lessons you teach. 
%\item Deviants: You don't have to be what they made you. 

\item Mad Scientists: So, if I get the institute to apologize for laughing at your research, will you destroy the killbots? No? Do you even know what you want?!
\item Leviathans: A society is its citizens, not its leader. If you don't believe me, just look at those monsters.

\item Hunters: Defending yourself is admirable, but fighting outside society can lead to dark places.
\item Mortals: Stop calling it cheesy when it's true. The power really is inside you if you just believe in yourselves. 
\end{itemize}

\section*{Concepts}

Literal princess, class president, owner of the local pub, priestess, member of the legislature, mentor, activist. 

\section*{Inspiration}
Hinamori Amu, Madoka Kaname

\end{multicols}\newpage\subchapquote{Mender}{\textit{Medicine looks simple: A sick person comes before you and leaves healthy. But healing is really about people. Most injuries and illnesses trace their roots to someone's decisions. Maybe they were careless and had an accident, maybe they were lazy and let their body weaken from a lack of exercise. Healing is complicated because you're not treating a body, you are always treating a person.}}
\phantomsection
\label{scha:Mender}\begin{multicols}{2}\raggedcolumns

\noindent All of the Nobility feel the pain of others, but Menders are those who can't bear to live with it. Their first impulse is always to treat pain, either physical or spiritual. These were the great sages and physicians of the Kingdom, who studied the science of the human conditions on levels undreamt by modern science. A good Mender is always looking to see how she can make those around her better. Today's dark world is like a ravenous beast, they say, and it mangles us all, all the time. Without care, none of us will ever be strong enough to break out of the mire that swallows us whole. Nobody should be made to lie crippled as their dreams run ahead of them.

Menders are typically described as either motherly or childlike – always blindly over-protective or naively concerned with the problems right before without seeing the big picture. They are natural caretakers. They easily take responsibility for those around them and are quick to worry. Often, they will display great maturity, although they are also frequently overcome by feelings of failure or helplessness when they lose a patient. A large number are pacifists, even in this dark world. Their uniting characteristics are a deep respect for life and horror in the face of suffering. The latter may sometimes be a weakness, but the best turn it into an iron resolve. 

\section*{Dreams of the Kingdom}

Menders' dreams are messianic in tone. They see themselves healing the sick and the lame with a touch or consoling those who cry and comforting those who mourn. In the hospitals and temples of the Kingdom, they nurse even the most desperate patients back from the edge. As the skies darken, though, all around them rot in their living graves and beg for death in a resounding chorus that deafens the would-be healers' very thoughts. 

\section*{Magic}

Menders have affinity for the Perfect, Restore, and Shape Charm families.

Starting Menders get one transformed Attribute dot in either Intelligence or Dexterity. 

\section*{Duties}

Menders regain magic whenever they provide consolation, comfort, and aid. They can volunteer their time as medical facilities or to provide home care, work as part of emergency hotlines, design new treatments, be a shoulder for someone beset by loss and grief, or do any other number of things to relieve the pain of others. 

A Mender can work with machines as well as people, a broken brake line can cause as much suffering as a broken rib. In the end though, it really does come back to suffering. A Mender builds and repairs machines to help people rather than for the simple joy of creation. 

\section*{Oaths}
\noindent \textbf{First Oath:} To be a healer is a sacred trust. My patients put their bodies or mind in my care, above all else I vow to do my patients no harm and respect their autonomy. The healer's sacred duties have survived the fall of our kingdom and I am proud to uphold their traditions. \\

\noindent \textbf{Second Oath:} When my help is needed, I will provide. Now, I don't have to be stupid about it -- if I don't have the skills it is not my help they need and I shouldn't make things worse by taking on a task I'm incapable of. If a dangerous creature seeks my aid, I can take precautions to protect myself and my friends even as I administer that aid, and I don't always necessarily need to give aid in the particular manner requested or unconditionally. And it's ok to prioritize or delegate; I can't be in two places at once. But I am here because people are hurting, and to ignore people's pain is to betray my Calling. \\

\noindent \textbf{Third Oath:} I am here to mend, not to break. I shall do no harm to another person. Even in cases of self defence or defending innocents, with my magic I can resolve the situation without hurting anybody. Of course, Darkspawn are no longer people.

\section*{Stereotypes}
\begin{itemize}
\item Champions: Bring the wounded back here, I'll be ready. 
\item Graces: Stop making speeches about the stars and look at how people live in the gutter. 
\item Seekers: The truth isn't hidden, it's right in front of us. People need our help. 
\item Troubadours: Just keep their spirits up for me. 
 
\item Vampires: Incurable. 
\item Werewolves: You could do so much good if you'd just let me try and duplicate your cellular regeneration. 
\item Mages: Give me one good reason you're not curing somebody's cancer right now, just one good reason. 
\item Prometheans: Bringing the dead back to life... growing replacement souls... it's possible? I need to know more. 
\item Changelings: Why do those feathers make me think of scars?
\item Sin-Eaters: I know you had one miraculous recovery, but as your doctor I must recommend you cut back on the drink, and the drugs, and the unprotected sex.
\item Mummies: I could help you remember yourself, why are you only interested in the wands I make? 
\item Beasts: I've fixed up enough of their students to know the lessons aren't worth the cost.
%\item Deviants: This should relieve some of the symptoms. I'm sorry, it's the best I can do. 

\item Mad Scientists: Do you want to talk about it? Of course I'm not calling you mad!
\item Leviathans: I can't decide if they need healing or cutting out.

\item Hunters: If you want, I'm willing to listen to why you're so hurt.
\item Mortals: No this isn't going to sting, I'm better than that. 
\end{itemize}

\section*{Concepts}
Emergency physician, IT support, miracle working technician, anti female genital mutilation or anti circumcision activist, hospice nurse, candy striper who's the hospital's real magic, doctor without borders, animal rescue, environmental clean-up technician. 

\section*{Inspiration}

Shamal

\end{multicols}\newpage\subchapquote{Seeker}{\textit{I followed my calling because I wanted to help people but once you learn to see you can't close your eyes. I've seen the victims I've helped turn around and abuse in turn. I've seen the heroes I looked to for strength crying in secrecy. But if you can see all that and still see hope you can do anything. The truth might be scary, but the truth will set you free.}}
\phantomsection
\label{scha:Seeker}\begin{multicols}{2}\raggedcolumns

\noindent The desire to know is universal human trait, embedded in the heart and soul and genetic code of every man and woman on the planet. What isn't common is the will to keep chasing after those answers when things start getting dangerous. That's what sets the Seekers apart. In their previous lives, Seekers were the detectives, scientists, and scholars of the Kingdom. Their dreams of the Kingdom are filled with colleges where the secrets of the universe are revealed and cities where no lie is uttered and not even the smallest misunderstanding is present to obscure the truth. When they dream of the Cataclysm, they see grand libraries and academies burned to the ground, and lies entwined with lies in a web of deception so deep that even the most fundamental truths are uncertain.

Seekers are a varied bunch. Many Seekers act like detectives or journalists, ferreting out crime and deception from the communities under their protection. Others resemble the scientist-heroes or adventurer-archaeologists of the pulp era, braving danger and using the secrets they discover to fight evil and better the world. Some Seekers act as stoic guardians of Things Man was Not Meant to Know, keeping forbidden knowledge out of unworthy hands. Others are stalwart enemies of lies and deceit, and those who would use them to pacify and control a population. But no matter how they approach their duty, all Seekers share the common goal of unearthing hidden secrets and more importantly, using those secrets to make the world a better place.

Seeker often prioritize mental attributes to aid them in their search for the truth. Resolve is prized by all Seekers to keep searching after the truth despite hardship or weariness. Wits, Larceny, and Investigation are prioritized by Seekers who favor classic detective work, along with Manipulation, Empathy, and Persuasion to tease the truth out of reluctant witnesses or suspects. Those who were scientists or scholars before Blossoming often possess high Intelligence. Subterfuge comes in handy when a Seeker needs to investigate things without creating suspicion.

\section*{Dreams of the Kingdom}

Seekers dream of secrets and discovery. They walk within libraries beyond imagination and delve deep into forgotten temples to uncover the lore of the past. With the knowledge they uncover the people build magnificent edifices and intricate devices that amaze and bring joy to people's lives. Then, a Seeker's dreams turn to the fall, the secrets she uncovers speak only of her own failings and the failings of that she holds dear. She desperately searches libraries for answers that do not exist while the darkness creeps ever closer. 

\section*{Magic}

Seekers have affinity for the Appear, Govern and Learn Charm families.

Starting Seekers get one transformed Attribute dot in either Intelligence or Wits. 

\section*{Duties}

A Seeker regains Wisps whenever she uncovers new facts or information, and whenever she teaches another something that they do not know. A Seeker could be a scientist or explorer uncovering secrets no man has ever known. She could be a detective uncovering that which unscrupulous individuals try hard to keep hidden. She could be a teacher passing what she knows to the next generation. Many Seekers do a both. A reporter uncovers a political scandal, then she publishes it. A Professor does exciting scientific research, then writes a paper and gives lectures on it. 

\section*{Oaths}
\noindent \textbf{First Oath:} My eyes shall always be open to the truth, I will never stop investigating for fear of learning that which I do not wish to know. \\
\noindent \textbf{Second Oath:} The knowledge I learn serves a higher purpose than satisfying my own curiosity, I will never hold back information that could aid society or prevent serious harm. If revealing that information also risks causing serious harm I will choose my response with care and discretion.  \\
\noindent \textbf{Third Oath:} The world is full of mysteries and they should not be left unexplored. I shall work to increase my understanding of the world every day.


\section*{Stereotypes}
\begin{itemize}

\item Champions: I understand how you feel, but punching people isn't the solution.
\item Graces: If only they asked any questions that actually mattered...
\item Menders: I admire your dedication, but you're treating the symptoms, not the disease.
\item Troubadours: Stop prattling on about art and actually do something.

\item Vampires: Friggin' lying bastards hurting people and lying all over the friggin' place. I hate them.
\item Werewolves: Great sources of lore if you convince them not to bite your head off.
\item Mages: A riddle wrapped in a mystery inside an enigma. Nothing that secretive can be any good. 
\item Prometheans: Don't blame them for what they are. Who's the bigger monster, the creature or Dr. Frankenstein?
\item Changelings: They're full of valuable insights. It's gaining their trust that's the chore.
\item Sin-Eaters: They probably have useful information to share. If only we could get them to talk about anything but doom and gloom.
\item Mummies: The truth is still there, burred beneath the sands. I could help you find it.
\item Beasts: Something isn't right here. If they were created to teach lessons why are they so bad at it?
%\item Deviants: Start from the beginning and tell me everything, the world needs to know. 

\item Mad Scientists: You'd be surprised how much they know, if they ever bothered to remember there's a difference between what they know is real, what they think might be real and what they think is cool.
\item Leviathans: There are some things even I don't want to know.

\item Hunters: They're great in a fight, but half the time they keep shooting up what I'm looking for!
\item Mortals: People have been lying to you your whole lives. I'm gonna put an end to that, just you wait and see. 

\end{itemize}

\section*{Concepts}
Research scientist, intrepid reporter, spiritual scholar, meddling kid, dreamlands explorer, hippie psychonaught, going to change the world from your garage, sociologist, Google-fu grand master, Nobel laureate.

\section*{Inspiration}
Yuki Nagato, Sakura Kinomoto, Ms. Frizzle, Yuuno Scrya 

\end{multicols}\newpage\subchapquote{Troubadour}{\textit{I spend my life under the spotlight. Every word I say and every song I sing is taken apart by people looking to make a quick buck or claim I endorse them. You have to keep dancing to stay one step ahead; that's how you say what words cannot express to tell the world what it needs to hear. Come on, it's showtime! }}
\phantomsection
\label{scha:Troubadour}\begin{multicols}{2}\raggedcolumns

\noindent Soul is not a silent, static egg that we must shelter from the outside world until the time comes for it to hatch. Troubadours know the truth – the soul is a living, breathing maelstrom and unless we open up and let it out, it can smother to death within us. And the Troubadours have the keys to let it out.

Each Troubadour is a master of some artifice – singing, dancing, painting, sculpting, baking, cracking jokes, or even just folding paper cranes. In their past lives, they were the queens of bards, storytellers and chanters, masters of divine instruments who imparted the wisdom of the ages in grand halls and shady grooves. Now their mission now is to set us free from ourselves. Art, the Troubadours teach, gives shape to things that people can't or won't express themselves. Every song or painting or performance lets the silent and unassuming realize and crystallize the beautiful internal forces they otherwise overlook. The people of the Fallen World suffer through with souls buried under the rubble of the Kingdom. The Troubadours are archaeologists of Joy.

Most people would describe Troubadours as people with their heads in the clouds. If they aren't visualizing a statue or planning a new routine or reviewing the forms of a plot, then a Troubadour is admiring the natural art of the world around her – the patterns woven out of everyday life and the people who populate it. Troubadours are vibrant, expression personalities, not shy about baring their souls. Their chosen art tends to consume their life and they want to involve as many people as they can in It. Unfortunately, they are artists first and foremost – they easily get lost in a world of symbols and frequently can't understand people's impatience with their “frivolous” hobbies. 

\section*{Dreams of the Kingdom}

Troubadours dream of their own artist ambitions displayed on an epic scale – murals that span cities, songs that resonate to the heavens, whole countries dancing. In these dreams, they are the conductors that make the whole world come to life with their energy and enthusiasm. The Cataclysm descends and they find themselves singers in the land of the deaf, painters in the land of the blind, clowns under the flag of melancholy, wandering a cold world as they starve in rags. 

\section*{Magic}
Troubadours have affinity for the Appear, Inspire and Shape Charm families.

Starting Troubadours get one transformed Attribute dot in either Presence or Dexterity. 

\section*{Duties}
Troubadours draw magic from the release of others' souls. Their duty encompasses providing art that inspires or consoles those who otherwise draw into themselves, encouraging and teaching those looking for a way to express themselves, and energizing those around them with their own vibrant talents. 
\section*{Oaths}
\noindent \textbf{First Oath:} I will use my talents to inspire, not to dissuade. My art should hearten my audience, never make them ignore and accept injustice. \\

\noindent \textbf{Second Oath:}My art will ever be for my audience. I must not become so wrapped up in my own ego that I forget the people around me. Without the mortals who need it, my art is merely a selfish exercise in solipsism. \\

\noindent \textbf{Third Oath:} My art will be. As long as I draw breath I shall not cease creating and sharing my art to improve the world.

\section*{Stereotypes}
\begin{itemize}
\item Champions: When you can use it to aim public opinion at a threat to the community, a song is a sharper weapon than a sword.
\item Graces: Laws don't make a society, art makes a society.
\item Menders: You can heal their body, I can enrich their souls.
\item Seekers: Do you remember your teacher's lecture on the states and their capitals? No? Do you remember that song from Animaniacs? Yeah, it's like that.
 
\item Vampires: You dress yourselves up and act so elegant, but you're really just wearing funeral makeup.
\item Werewolves: Such strong emotions, hold still. I need to get this on canvas. 
\item Mages: I don't understand, if you experienced a revelation like that how can you not be shouting it from the rooftops. 
\item Prometheans: To literally touch god through your creations...
\item Changelings: They're very secretive, but you can see their pain in their art if you know how to look.
\item Sin-Eaters: Such a dynamic and energetic people, pity there's no substance behind it all.
\item Mummies: I admit it, I'm jealous. My works haven't stood for thousands of years. 
\item Beasts:  I used to tell my little sister monster stories when she was naughty. I still have the books, they never knifed anyone and stole their wallet.


\item Mad Scientists: I know how it feels to be inspired, but there's really no excuse for what you did.
\item Leviathans: That's just horrible, and maybe just a little bit beautiful. Wait! What am I saying?!

\item Hunters: You fight monsters? Without any special powers to help you out? And you win? [bites lip] Do you want my number?
\item Mortals: I've got something to show you, I think you'll quite like it. 
\end{itemize}

\section*{Concepts}
Painter, baker, pop star, surrealist who's works contain hidden wisdom, dancer, star athlete, graffiti artist, practitioner of traditional styles. 

\section*{Inspiration}

Creamy Mami, Fancy Lala, Lucia Nanami, Kongou Shin, Haruhara Haruko, Jem and the Holograms, the EBA 



\end{multicols}\newpage\subchapquote{The Queen of Clubs}{
\noindent\textbf{AKA:} The Matron of the Forests, The Mother\\
\noindent\textbf{Followers' Epithets:} Wilds, Turtles\\
\noindent\textbf{Kingdom:} Wen-Mung}
\phantomsection
\label{scha:Clubs}\begin{multicols}{2}\raggedcolumns

\noindent \textit{In any battle your greatest foe is yourself. You must ask yourself why you fight, what can you hope to gain and is it worth the risk of what you might lose? The enemy who says you cannot back down is none other than yourself; overcome it and look at the other options. You could surrender, often it costs you nothing but an apology and some pride. You could examine yourself, what did you do that lead up to this moment. Why couldn't you live here in peace? If you can learn to live in harmony with your surroundings, isn't that a better result than whatever you could win in a fight?}

\textit{You can't fix the world by yourself, but the best place to start is within.}\\

\noindent The Queen of Clubs leads those Princesses who seek balance and harmony, they believe that a better world can be achieved if we just learned to live peacefully with each other, with the world around us, and most importantly: with ourselves. 

\section*{Tales of the Kingdom}
The Queen of Clubs has no palace and wears no crown. Her court travels at a slow walking pace through the forests of Wen-Mung. Wherever they rest their tired feet the trees grow to shelter her from the elements. With nothing but bare soil for a throne The Matron of the Forests welcomes her subjects. Of all the Queens she is the most intimately connected with the eternal cycle of life, death and rebirth. She allows herself to age naturally, and in her current iteration of the cycle has grown into a slightly plump middle aged woman who radiates warmth and kindness. She wears only spare clothes gifted to her by the working women of her kingdom and rejects all finery. When she speaks her words are plain but carry the wisdom and empathy of lifetimes.

A good way to describe a conversation with the Queen of Clubs is patient; protracted works too. The Matron of the Forests sees no need to hurry. She lets the Princess speak at her own pace and get to the end of her story before she replies. When presented with a tough problem she likes to work up to it, focusing first on a minor but easily solved parts of the issue. This usually includes coaching the Princess about the right attitude to a problem before fixing the problem itself; though it must be said that the Matron of the Forests sees the first as a vital part of the second. 

\section*{Philosophy}
The Queen of Clubs asks a question: \textquotedblleft The world speaks to us. Have you ever stopped to listen?\textquotedblright

\subsection*{Through Harmony; Light}

The world is a place of conflict. Not just the obvious conflicts where nation fights nation but also the conflicts that arise from the simple act of existing. The Queen of Clubs teaches her daughters that they and everything in existence has a nature. It is from the struggle to be true to this nature within the confines life imposes that suffering arises. Through introspection and self discovery a Princess can find a way for her nature to exist in harmony with the world and free herself from suffering.

\subsection*{Things Change My Dear}

Harmony is not a destination, it is a way of moving with the world. Everything has a nature and change is part of that nature: Water flows, plants grow and planets dance around the stars. With great effort a Princess can maintain a harmony by opposing negative change but this is a false path. The true harmony is effortless. A state of grace where every action moves with the world and every change creates a new harmony free from discord. When the world changes the court of Clubs must move with it, and they shall becomes greater from the experience.

\subsection*{Be The World You Wish For}

To the Queen of Clubs the world is the lives and actions of every living being. From this simple truth she forms her final philosophy. Harmony cannot be won through conflict for the very act that claims to create a world free from conflict becomes part of the world. To create a better world a Princess must become a teacher and a mediator who lives a life of harmony and helps others achieve the same through her example and guidance. 

\section*{Duties}

A Turtle's Duties are to bring things together. People certainly but that's only a small part of it: Man and nature, reason and emotion, city and countryside. The Wilds dream of a world where what were once opposing opposites are now part of a balanced whole. 

Champions of Clubs help people directly. They are park rangers and wilderness guides who help people brave the dangers of the natural world so that they may see it's beauty first hand. More than any other Calling it is the Champions of Clubs who earn the epithet of Turtles. In battle they favour defensive tactics and seek to gradually wear their opponents down with a minimum of harm. 

Many Graces are environmentalists who seek to preserve natural world and preach the importance of living in balance with our surroundings. Other Graces are mediators without peer who work to bring conflicting parties together in a harmonious balance. The lessons of harmony are not limited to those with an external enemy. The Queen of Clubs teaches that harmony begins with inner balance and many of her Graces are teachers who help their students learn to live in harmony with themselves and the world around them. 

In the Court of Clubs Menders are rarely simple healers. Under the Queen's philosophy the disease itself is often a symptom and her Menders teach their patients to live a lifestyle that avoids further sickness. The court's focus on the natural world attracts membership from vets and environmentalists who work to clean pollution. Partnerships between Menders and Graces are common, the Mender deals with the physical symptoms and the Grace teaches a way of life that removes the root cause. 

Seekers delve deep into the jungle seeking new herbs, undiscovered animals or just beautiful vistas. Others travel beyond the borders of earth. They delve into the lore of the Spirit wilds, the hidden half of the natural world, and try and use the wisdom they find to shape and heal the spirit community around them. A few Seekers blur the lines between Clubs and Diamonds as they work on renewable energy and cleaner technology. 

Troubadours are inspired by nature and paint magnificent scenes of the natural world, others work with nature directly as landscapers, flower arrangers or bonsai artists. A surprising number are architects and urban planners who attempt to reintegrate nature with cities in the most literal manner.

\section*{Background}

Appropriately for a Court so focused on Harmony, the Wilds are often classified into two groups with little conflict or divisions between them. On the one hand you have the naturalists: Everyone from professional park rangers and conservationists to people who really love their pets. They are drawn by Clubs' belief in harmony with the natural world and often leap for joy when they learn how Legno lets them talk to animals. 

On the second hand you have the spiritualists: From gurus and philosophers to pacifistic hippies and suburbanites yearning for something with more meaning than materialism. That is not to say that every spiritual Princess is drawn to the Wilds. The Court of Clubs has its own philosophies which draws students interested in learning about harmony, balance and the quest for inner peace.

When the Court of Clubs gathers it is not rare for a laboratory biologist to talk long into the night with a self taught student of Taoism, each delighting seeing the world through eyes that on the first glance are so different from their own. The Queen herself knows that the two sides are both the inevitable result of each other and smiles every time one of her sons or daughters experiences the joy of learning so for themselves. 

\section*{Signature Emotions: Harmony and Tranquilly}

The Queen of Clubs teaches that harmony is the goal. To be in harmony with your community, your surroundings, and yourself is to be free from suffering. Everybody deserves nothing less. And if she is to achieve harmony a Princess must first embrace tranquillity. Through calm and inner peace a Princess can accept the world around her with all it's flaws and love it for what it is. It is through accepting the world's imperfections that she can see how to weave a harmonious whole out of imperfect parts.

\section*{Character Creation}

Composure is highly desired among the Turtles to examine their own feelings while it takes considerable Resolve to always keep striving for self improvement. Wilds always put at least some effort into keeping themselves fit, meaning that Physical Attributes are at least average. Most Wilds make a point of familiarizing themselves with nature, so skills such as Survival and Animal Ken are common. Athletics come in handy during hiking, swimming, climbing, and otherwise navigating around in the wilderness (or at least the state park). Wilds who wish to convince others to embrace the natural world often pick up Persuasion and Expression, while those who have a theoretical understanding of nature often possess Science. The spiritual mindset of the court lends itself well to Academics, Wilds interested in Spirits often pick up Occult.

Socially the Court of Clubs tends to have fewer dots in Allies than Contacts. It is the way of Clubs to coexist, not control. Circle however is often staggering as the Court's focus on co-existence brings them closer to whoever is around them. It is not unknown for a Wild to count an entire community as their Circle, whether that community is an ordinary suburban neighbourhood or a commune built on principles of friendship and love seems to make little difference. The outdoor life enjoyed by many Wilds makes Physical Merits common, many are exceptionally healthy people who enjoy a Natural Immunity and a good Direction Sense is common among both wilderness explorers and hikers. Mental Merits are rarer, but a Meditative Mind and Holistic Awareness are both common in the Court of Clubs. In conversation many favour the Wisdom style.  

\subsection*{Heraldry}

In their Transformed state, Wilds tend to be more subdued than other Princesses. White, green, and blue are common colours, although occasionally red, orange and yellow make appearances for those with an autumn theme. Ivy and leaf motifs are also common, with water themes being only slightly less popular. Outfits tend toward the elegant and practical, the Wilds favour tough clothes well suited for the outdoors, but also with a sense of balance that lends grace to their appearance.

\subsection*{Echo}

A Wilds echo tends to give people feeling of serenity. People feel calm and peaceful around the Princess, as if her mere presence somehow brings the world into greater harmony.

\subsection*{Practical Magic}

The Queen Mother teaches her students that conflict is a cycle that must be broken and her Court can find strength or victory through refusing battle. Once a scene, when a Wild is the target of a resisted or contested action, she may spend 1 Wisp to twist her opponent’s action in her favour. Her resistance doesn’t apply, or she gets no successes in the contest; but if her opponent succeeds she imposes a negative Condition on her opponent or takes a positive Condition herself. The Condition can be any Condition that flows naturally from the situation, if none seem suitable an improvised condition [CofD 76] can be used. As a rule of thumb the Condition should be of a similar strength to the effect that the other character is trying to inflict on the Princess. The Storyteller remains the final arbitrator of which Conditions are acceptable.

High Belief softens the blow from the force the Wild redirects. At Belief 8 she adds 1 to the successes her opponent needs for an exceptional success; at Belief 9 she adds 2, and at Belief 10 she adds 3. So at Belief 8 her opponent needs 6 Successes to score an Exceptional Success. 

\subsection*{Invocation: Legno}

The Queen of Clubs teaches the principles of harmony and inner strength in the form of Legno and the Turtles find it easier to master than any other. Legno encourages a Princess to work with, not against. To find balance not overpower. To defend not attack. It also guides a Princess to seek inspiration in the natural world.

Legno applies at no cost when the target of a Princess' Charm is a living plant or wood, and when her Charm's target is a natural, non-sapient animal. It also applies without cost when a Princess intends to restore things to their natural state, to restrain the reckless and calm the fearful, to find a harmonious compromise between two positions, to protect against an attack, to prevent or calm violence without committing violence herself, to conserve and cultivate over time, and to influence without drawing attention to herself.

Those who would make use of Legno must live for harmony in all things. A Princess who attempts violence against someone not involved in combat, or who helps another to do so, may not apply the Invocation until the next sunrise. Self-defense, protecting a non-combatant, and using Charms to stop a fight without violence are unexceptionable, but if a Princess gives tactical support to an ally and the ally attacks one who has not begun to fight, she loses access to Legno just as if she had struck the blow herself. Darkspawn and other creatures of the Darkness (except Darkened) are not covered by this ban. 

\section*{Quote}

You can't control the world, but you don't need to.

\section*{Stereotypes}

\begin{itemize}
\item Diamonds: They seek knowledge only for the sake of knowledge. They build a future with no idea where it leads. Yet we must ask ourselves: Would we have time to contemplate without their agriculture? Would we understand the life of the butterfly without their tables of numbers? Ask them to free us from necessity, our work comes after.
\item Hearts: They are living in traditions built for the world of yesteryear. Either they can hold the world back or the world can trample them underfoot. It cannot end well.
\item Spades: Flexible, adaptable they can fit anywhere yet they lack the permanence to ask the same respect of the world. Somehow I don't think that matters to them.
\item Swords: If you are strong enough, you can move the world, but unless you teach the world to move itself it will always return when you remove your hands.

\item Tears: The world changed, until you can change with it you will suffer. I would like to help you, but only you have the power to help yourself. 
\item Storms: When you learn to take every blow the Darkness throws at you, then you will learn to fight without falling to hate. 
\item Mirrors: You claim that you shall solve every problem yourself by the virtue of your crown, but your claims of strength are to hide from a weakness within. It is an appropriate paradox, in a mirror everything is back to front.       

\item Vampires: We’ve all heard the tales, that they are creatures of pure Darkness who run away from the Light’s touch. I didn’t want to believe them. Then I met one.
\item Werewolves: Their mission is to keep two worlds apart, ours is to bring two worlds together. But I wouldn't recommend telling them that. 
\item Mages: I've been watching you, I saw you drawing maps but did you notice every path you found was first found by the ants? This glade that delights you so, did you see that every moss and mushroom grows to face the centre? You can feel the Earth's power, but you cannot master it until you listen to her heart. I could teach you so much if you only had the humility to listen.
\item Prometheans: I met one once. It’s been the only time in my life were I was sure I had met something that couldn’t belong in an harmonious world. The next time I met him, he was just a man. How could anyone overcome such a burden? I have to know.
\item Changelings: You are trying to balance two natures in your heart. I would say we were alike if your other nature was not so different. Instead I offer the kiss of a Princess for luck.
\item Sin-Eaters: Behind their denial of death is a clear understanding that death is a natural part of life. 
\item Mummies: I can think of no curse greater than an unchanging immortality.  
\item Beasts: They adapt to new environments by finding justifications to avoid change. Is it any wonder why they bring harm to all they touch?
%\item Deviants: If you can't go back to who you were, go forwards.
\item Mad-Scientists: Trying to force the world to change is bad enough. Trying to force the world to change just because you can is madness. I'd offer guidance if there wasn't a forcefield in the way.
\item Leviathans: You look old, timeless, like something from the early days when the world hadn't learned there was more to life than sex and violence. I'm sorry, but we've come so far since then and until you do the same there's just no place for you here any more.

\item Hunters: The only thing you can say they have in common, is that they define themselves by what they are in opposition to. They do good for the world, but they could do so much more if they found themselves first.
\item Mortals: We live the world we want for them so it may come to be. It’s not easy, but I believe they are worth it. \end{itemize}

\section*{Inspiration}

Nausica\"{a}

\end{multicols}\subchapquote{The Queen of Diamonds}{
\noindent\textbf{AKA:} The Crystal-eyed One, Lady of Clear Water, The Philosopher Queen\\
\noindent\textbf{Followers' Epithets:} Crystals, Lights, Hope-Engineers\\
\noindent\textbf{Kingdom:} The Danann Archipelago}
\phantomsection
\label{scha:Diamonds}\begin{multicols}{2}\raggedcolumns

\noindent\textit{Contrary to popular belief, people aren't stupid. It's just that people don't think, or they don't know how to think, or they just didn't have the right facts. If you don't think, you miss all the flaws in your plans. If you don't have all the facts, you make stupid assumptions because you didn't know any better. If you don't think and don't understand the world around you, you'll put all your efforts into a plan that simply won't work. And then things go wrong. Sometimes, they go very, very wrong.}

\textit{That's why thinking and knowledge and education are so important. If you know why a problem happens, you're halfway to solving it. If you think clearly, you can find that solution. And if you educate people, you give them the ability to solve their own problems. By understanding the world, and teaching others to do the same, we can find a way to make it brighter.}

\textit{So. Want to help build a better future?}\\

\noindent The Queen of Diamonds teaches that the Light of Hope is best accompanied by the twin lights of knowledge and reason. The Hope-Engineers say that to create a better world they must harness their passions and emotions through careful planning and forethought. A better world will be built, but first it must be designed. 

\section*{Tales of the Kingdom}

From her palatal starscraper the Queen of Diamonds rules the Danann Archipelago. Upon a throne of whale ivory and pearl sits The Philosopher Queen: Regal as a glacier, six foot tall and adorned in the latest sleek future-chic fashions of her court. Here she judges the petitions and arguments of her subjects. Her replies are lengthy and explain the facts, theories and reasoning behind every decision. Every word spoken by The Lady of Clear Water is precise and rehearsed, as calm as a moonlit pond.

Those who know the Queen in private see a very different person. Barely reaching five foot stands a young woman just entering her twenties, wearing casual clothes, unkempt mousy hair and a tee shirt adorned with an obscure and highly technical joke. She even wears glasses having never bothered to fix her eyesight. Far from the regal Philosopher Queen she stumbles over her words, gushing like a bubbling brook in her excitement over the latest technical and scientific discoveries. When gifted with some unseen device or research from Earth she squees, yes squees, in delight. Among her many friends and frequent one night consorts she values nothing more than someone with whom she can share her boundless passion for learning.

The Lady of Clear Water tends to assume that the sort of Princess who pledges to the Court of Diamonds is the kind of person she can introduce to her private self; she is usually right. When talking to the Queen in her public persona it is best to speak as she does, clear, precise, with a focus on the facts and an explanation of your reasoning. In private talking to the Queen of Diamonds is like sailing down a river; her excitement leads the conversation and the Princess must choose whether to follow or fight the currents to steer the conversation in more practical directions. Some Princesses become tributaries, their own excitement joining with the Queen's to become greater than either alone; it is these Princesses most likely to win the Queen's friendship and the royal favour.

Wherever the Queen travels Plato, her ever-loyal owl, is sure to follow. To her court, and even her close friends Plato is nothing more than her pet and her idea sounding board. In truth Plato is the Queen's loyal spymaster. Among the Court the wise owl is second only in intellectual might to the Queen herself and among the strongest magical practitioners.

\section*{Philosophy}

The Queen of Diamonds has taken as her motto, \textquotedblleft Think, then act.\textquotedblright

\subsection*{Through Understanding, Triumph}
Everything has a cause; everything has a purpose; nothing that happens is inexplicable. The servants of the Philosopher Queen take the duty of finding causes, purposes and explanations, and bear it with pride. The Crystals would inquire, as a rule, for the sheer love of knowledge even if no benefit could be expected. Nonetheless there usually is some practical use to the Lights' studies. When they come across a problem, the Hope-Engineers analyse it carefully, and reach as complete an understanding of it as they can; then they strike, they cut the problem off at the source and watch it drain away. Analysis is obligatory, assumptions and intuitions can be a starting point but they must be tested, since an ill-planned strategy is usually worse than none. When no problem is at hand a Light studies new frontiers to push the boundaries of knowledge or re-examines old traditions and accepted wisdom, looking for errors or outdated ideas. The more widely a Crystal has studied, the more swiftly she takes hold of a new 
problem and supplies its correct solution.

\subsection*{Guide the River's Path}
A Princess can solve many problems with magic or the raw intellect of her perfected self, but even in the days of the kingdom Noble magic was finite resource. The Philosopher Queen teaches that the best solutions are those that sustain themselves. It is for this reason (among others) that the Hope-Engineers are enamoured with mortal technology. Instead of curing diseases with magic, create vaccines and medicines to banish the disease forever. But at heart any self sustaining answer will find the Queen's favour. While a Crystal could personally organise an economy, it is better to design an infrastructure that makes the economy easy for anyone to organise. Whatever problem she faces, a Crystal should look at the big picture and craft an elegant solution that will last long after her reign has ended.

\subsection*{Light Shines Through Clear Minds}
However unusual a fact is, however much a theory encompasses, it has no power if people know nothing of it, or lack the training to understand it. So, the Hope-Engineers take the further duty of education, and of clear communication. Knowledge and knowing how to reason are invariably good, lack of either a weakness, withholding of them a dereliction of duties. A Crystal pushes everyone she has any influence over to learn, to think, to question, and to teach, and she pushes herself most of all. To the Lights the bond between teacher and student is sacred, and must never be tainted with lies or laziness. 

\section*{Duties}
The Lights' Duty is to use her brain. The Courtiers of Diamonds seek positions where problems can be solved with thought and intelligence alone. Sometimes it falls upon them to personally implement the solutions they conceive. Sometimes society expects them to start at the bottom and learn the ropes. But whatever situation she finds herself in, the ideal is to earn a Mandate to use her mind.

What Champions of Diamonds lack in combat power they make up for in strategic ability. As much a general as a fighter the Champions of Diamonds have the mind to plan a grand military campaign and organise the logistics of keeping her army fed and battle ready. When battle is joined her quick wits let her adapt to changing circumstances and match her foes gambit for gambit. If so many Champions of Diamonds are seen with shield and spear in hand against the Darkness, that's because so few of them have an army to lead. While her Nakama will benefit from a well thought out plan to raid a Darkened cult, a group of five can scarcely afford to leave their Champion away from the front line.

The Graces in the Diamond Court might not be the best at managing people but they are the best at managing their kingdoms. The Grace's pen closes loopholes in the tax code, defines laws in clear and understandable terms or just balances the national budget. So long as she can keep her head above the confusing tangle of special interests and the reins of policy away from those who pander to badly thought out populism she should have no trouble running a kingdom.

Seekers of Diamonds are perhaps the purest expression of Diamond's Philosophy. They truly believe that knowledge is inherently good and only application can be evil. Deep within their ivory towers the Seekers of Diamonds study everything, from atomic physics to the Twilight Invocations. The future is built on what they discover.

If The Queen of Diamonds' Seekers are scientists, then her Menders are the engineers. A Mender of Diamonds is more likely to be found in a laboratory researching new cures and vaccines than in the hospital, helping patient after patient. Rather than fixing a broken machine, she will design a new machine free of the fault that broke the first model. The Hope Engineers can be found in research and development teams, developing technological policy, contributing to open source, and running crisis response centres.

The Troubadours of Diamonds tend to dedicate their talents in a very specific direction. Their duty is simple: make learning fun! Whether they write a sci-fi epics where the plot hinges on an accurate understanding of science or whether they're the cool chemistry teacher who manages to tie every topic into a practical that ends with something blowing up. When Troubadours of Diamonds are teaching, you might not even know you're being educated.

\section*{Background}

No matter what cultural or geographical background they come from, the court of Diamonds were always people who appreciated the importance of thinking things through. Access to education is not nearly as important as some might think, but whether she was educated in a first world university or self-taught in a third world slum a would be Light tends to make the most of whatever educational opportunities she had. The Hope-Engineers recruit heavily from the ranks of teacher's pets, library bookworms and self professed nerds. Professors and PhD students are rarer, most Princesses Blossom before graduation. 

Despite their Queen's love of science and technology, it is not necessary for a Hope-Engineer to be (or aspire to be) a scientist or work in a technical field. The Court's welcomes anyone who believes in the value of thinking things through. Only around half the (adult) Court can be said to work directly in an academic or technical profession, which includes a large crop of librarians, teachers and other educators. Regardless of background, most of the Court sincerely believes in the potential for technology to create a better world. 

\section*{Signature Emotions: Curiosity and Wonder}

This world is amazing. Science, nature, history, culture, there is more to learn than can ever be known. Let others look at the world and despair, Diamonds see beauty and wonder everywhere they look; and in that beauty they see a path to a better future. In the shell of a darkling beetle is the secret to bringing water to the desert. In journals written and read by obscure mathematicians are formulas that hold the key to a fairer and more democratic electoral system. Everywhere a Light looks she sees things worth fighting for and ways to win that fight. 

And if the world is so full of wonders, how could you not be curious? How could you not want to see what lies around the next bend, under the next rock, over the next page? A crystal's curiosity compels her to always look further, look deeper. The more she learns the more problem's she'll know how to solve, and the more beauty she sees the more she'll want to learn. It's a cycle as endless as the tides. 

\section*{Character Creation}
Intelligence and Wits are both highly prized among Crystals. Intelligence is vital for making sense of data, while Wits is valued for spotting patterns and noticing important details. Many also focus on Composure or Resolve, either to keep their heads clear and emotions under control, or to keep one's mind on the facts no matter how strong their emotions run. If a Light doesn't have Investigation, she tends to pick it up fairly quickly in their efforts to get the facts of whatever problem she's facing. Many Crystals have at least a few dots spread across Computer, Crafts, Academics, and Science. Those who seek to understand the strange things of the World of Darkness often take dots in Occult. Persuasion and Expression also are fairly common, Persuasion for convincing other to follow the plan a Light has made, and Expression for explaining and teaching. Larcanay and Stealth see use in getting information that can't be obtained any other way.

Socially the Court of Diamonds favours Contacts, for the access to information it provides. Many find themselves joining scientific or engineering projects and a tight community with shared goals can easily lead to both Mandate and Allies. Those who act as ambassadors for their work typically pick up the Scientific Rhetoric conversation style. More than any other Court the Hope-Engineers enjoy access to mental Merits. Many have strange abilities such as an Edetic Memory or Common Sense granted to them by their Blossoming, others have studied all their life and have abilities like Encyclopaedic Knowledge to prove it. Champions of Diamonds usually have know Small Unit Tactics (the only reason they don't have large unit tactics is because it's so hard to get a large unit).

\subsection*{Heraldry}
Lights often have a scientific, technological, or futuristic feel to their heraldry. For some, their Regalia calls to mind science fiction: computer readouts built into sleeves or gloves, gadgets such as scanners, energy weapons, or palmtop computers. Many Regalia shed light in some manner. Clothing can be made from metal, plastic, or some strange material that defies categorization. The superhero look is fairly common among the Hope-Engineers, with colours tending towards white or the brighter shades of blue. Others have heraldry that call to mind professions of learning; many Lights wear glasses as part of their transformed identity or their Regalia, whether or not they normally need them. Some have long coats reminiscent of a lab coat, while others have vests, boots, and practical looking clothing suitable for explorers, archaeologists, or others who look to learn things far from a lab. No matter the style, the appearance of many Crystals features a distinct lack of anything that would be inconvenient or 
unwieldy a normal costume, despite the fact 
Princesses don't suffer any drawback from impractical garments.

\subsection*{Echo}
A Light's Echo is like a cold splash of water, it removes the sleepiness from the corner of your eyes and makes it easier to think.

\subsection*{Practical Magic}

The Queen of Diamonds' first lesson is \textquotedblleft think, then act\textquotedblright\, and her Princesses are her PhD students. Through intellect and magic, a Hope Engineer's planning is a thing of beauty to behold. The player declares broad goal like \textquotedblleft prove that GlobaBank is defrauding inner city customers\textquotedblright\, and the Princess spends at least an hour researching and planning. 

At the start of any later scene that fits her goal the player may declare that the character planned out this scene. Then for the remainder of the scene the Princess may, up to Acqua times, spend a Wisp and declare that she's planned ahead and her preparations are about to pay off. For example, if the players discover a finger print lock, she may declare that she already acquired a copy of the prints. The storyteller may call for an appropriate roll depending on the difficulty of the actions, and may disallow anything that would be an entire scene in of itself. Following someone and stealing a cup with their prints from a caf\'{e} could be done with a couple of rolls. Breaking into a CEO's secure mansion and stealing their prints would probably require a scene. The Princess does not need to make these rolls herself, she could delegate the actual theft to the Nakama's Knave. 

The Princess may not undo anything that's already happened in the scene. If the players trigger the alarm she cannot declare that she disabled the alarms, but she could declare that she hacked the telephone line so that the alarm makes a lot of noise inside the building but won't call the police (at least, not until the security team notice the problem).

High Belief increases the number of time a Princess may declare her planning bears fruit. She may do so an additional time at Belief eight, nine and ten for a maximum of three extra uses. 

Once the scene ends the Princess needs another planning session before she can use this power again. 
\subsection*{Invocation: Acqua}
The Philosopher Queen focuses the lights of knowledge and reason into the Invocation of Acqua, and the Crystals gain mastery of that Invocation more swiftly than any other. It is linked with the virtues the Crystal-eyed holds dear, such as clarity of mind, honesty in word and deed, deep understanding, and to the emotions of curiosity and wonder. It is also tied to the classical element of water, and to liquids in general. The invoked Charms and upgrades based on Acqua relate to these things and to phenomena connected to these.

Acqua applies at no cost when the target of a Princess' Charm is water or ice, or when she has spent at least one turn studying her target during the scene and nothing significant has happened to it since. It also applies without cost when a Princess intends to teach someone, to master a body of knowledge, to prove or refute a theory, to discover things previously unknown to anyone, and to carry out a strategy she has previously formulated.

The Invocation aids only those who are open and candid in all their dealings. A Princess who deceives another person in any fashion, be it simple falsehood, sophistry, or material omission, cannot apply Acqua for the rest of the scene, unless the deception is necessary to shield an innocent from harm, to keep a secret given her in confidence, or to protect a Princess' secret identity. A Princess may also refuse to answer a question and is not obliged to go into exhausting detail.

\section*{Quote}

Your heart can show you where you need to go, but only your head can get you there. 

\section*{Stereotypes}

\begin{itemize}
\item Clubs: You've got some good points, but I have absolutely no idea what you're trying to accomplish.
\item Hearts: We can learn a lot from the past, but you can't just imitate old ways blindly and expect them to work today.
\item Spades: They're excellent when it comes to making you think ... but are they \textit{never} serious?
\item Swords: When you need a grand, dramatic gesture, call for them. When you need subtlety, call for someone - anyone - else.
\item Tears: Look, it's obvious. If you didn't live in that dank hole you wouldn't need to steal the world's hope. Why don't you move out?
\item Storms: The ones who let their hate and anger control them, I can deal with. They ones who are still rational, I understand and that scares me like nothing else.
\item Mirrors: Someone needs to deflate their egos before they get someone killed.

\item Vampires: You're a problem. I solve problems.
\item Werewolves: I'll stay out of your way, try and stay out of mine.
\item Mages: They're dangerous \textendash\, not because they have fireballs, but because their understanding of reality is fundamentally wrong.
\item Prometheans: What \textquotedblleft are\textquotedblright\, you?
\item Changelings: The more I learn, the less sense they make. Thankfully, they don't seem to be an issue.
\item Sin-Eaters: They know a lot about working with the dead, but they don't seem to understand how death works.
\item Mummies: It sounds like your magic was far better understood than ours. Don't let that be lost to history, teach me. 
\item Beasts: One of us has studied psychology, educational theory and gets informed consent. The other needs to turn themselves in to the police before I shoot them.
%\item Deviants: See. This is why laboratories are supposed to have an ethics board. 
\item Mad Scientists: They're dangerous \textendash\, not because they have death rays, but because their understanding of reality is fundamentally wrong.
\item Leviathans: Find wonders in the deeps, whoever said that can't have known what's really in the deeps.

\item Hunters: They do study mysteries, I'll give them that. I wish they weren't so quick to shoot them.
\item Mortals: If we can teach them to think, they will make the world perfect.
\end{itemize}

\section*{Inspiration}
Jonas Salk, Sailor Mercury, Washu Hakubi, Watanabe Eriko, Hermione Granger

\end{multicols}\subchapquote{The Queen of Hearts}{
\noindent\textbf{AKA:} The Rose Bride, The Gentle Patroness, The Queen Regnant\\
\noindent\textbf{Followers' Epithets:} Jewels, Flowers (female), Gallants (male), Rocks (often derogatory) \\
\noindent\textbf{Kingdom:} Andarta\\}
\phantomsection
\label{scha:Hearts}\begin{multicols}{2}\raggedcolumns

\noindent\textit{I know what you think of society, you see nothing but the people at the top having dinner with politicaians and hiring expensive lawyers. A big system to keep the rich rich, and the poor poor. Read a history book some day. You will see that opression used to be the sword, not the law. That is why I believe in our society, where you can only see how far we have left to go I see how far we've come. We are on the right path, the future will be built upon the foundations of our history.}

\textit{So, can I count on your vote?}\\

\noindent The Queen of Hearts bids her subjects to restore the Kingdom's spirit to the modern world. The followers of Hearts believe a better world has already been built on the back of principles of civilisation and by serving the people as judges, leaders and lawmakers the Flowers believe they can make it better yet. 

\section*{Tales of the Kingdom}

The Queen of Hearts rules the city states of the Andarta Plains from her palatial estate. Within the Queen presides over her court from a throne of red granite and gold. She is everything a Queen could be. Regal, dignified, just, kind, possessed of a gentle sense of humour and easy upon the eye. Her skin is delicately adorned with ointments and perfume. She is clad in elegant and extravagant styles drawn by the royal tailors from all across the world and history to show the reach of their lieges influence. However she presents herself she is as composed as a flawless diamond, every word and action is deliberate. When she hears the request of her petitioners her court fills with the rustle of paper as an army of clerks search through scrolls and tomes for precedent of similar problems long since solved.

It is easy for a Princess of Hearts to talk with her Queen. She can learn the protocols in advance and use their predictability for reassurance. She will call the Queen by her title, and the Queen will call the Princess with her title. It helps that the Queen of Hearts can spot nervousness or awe a mile away and deftly puts the Princess at her ease. When the Princess speaks the Queen is an active listener, guiding the Princess with short questions but never taking the conversation until it is her time. In turn she expects the same courtesy from the Princess. While the formality and protocol can make it hard for a Princess to connect with the women beneath the crown; as two crowns talking for the good of their peoples it could not be more ideal.

The Queen of Hearts is a Queen wearing a mask. She is not a woman wearing the mask of a Queen, the few people who have seen her in private all agree she is a queen through and through. No, the The Gentle Patroness is a Queen with a crown for every occasion. In times of peace she is a public speaker without peer and a patron of charities and the arts. In times of uncertainty she is a diplomat, a peacemaker and a rock of stability. In times of war she is a warrior-queen, the first to charge and the last to retreat.

Among her titles is The Rose Bride. How the Queen of Hearts gave her own heart to another is a tale sung by bards throughout her kingdom. In her very first life, many many incarnations before she became a Queen, she married a young Prince. Their devotion to each other was so strong that in every life that followed the two fell in love and wed anew. Not once in any of their lives has either been unfaithful to their marriage. Though a diplomat by nature and profession the Prince Consort is also known as the Queen's Champion. He does not defend his' wife's life, for any foe that can threaten a Queen has nothing to fear from a Prince. He defends the Queen's honour, for in the court of Hearts it is considered proper that the stronger fighter in any duel takes a handicap. When one is as powerful as a Queen the only possible handicap is to send a champion to fight in your place.

\section*{Philosophy}

At the core of the Queen of Hearts' creed is potential for of society itself to be a force for good. She teaches her court how to unlock that potential and to fulfil their own responsibilities as part of society. All this rests on a simple foundation stone: \textquotedblleft A crown is a symbol that must be honoured\textquotedblright .

\subsection*{Authority Must Be Earned}

A Princess of Hearts, more than any other Noble, is expected to achieve a position of authority then use it to bring light to the world. In doing so she becomes the keeper of many people's trust; therefore she must be worthy of it, both in her own conduct and in the behaviour she tolerates. A Princess of Hearts \textit{must} earn authority and she must \textit{earn} authority. When a threat to a whole society appears, Flowers are the ones who rouse its members to action, coordinate their responses, and ensure everyone does what's needed. In ordinary times, Gallants are courteous, friendly and compassionate to all, and perform their duties with dignity and diligence; but they have their burden, heavier than most, and don't allow themselves to forget it. Moreover, they must not take their position for granted. If her society rejects a Flower's rule, or deposes her from it, the fault is hers not theirs; and if they proceed to disaster that she would have avoided, her failure to gain their trust is all the more grievous.

\subsection*{Flourish in Community}

No one can stand alone against all the Dark and survive; but people who trust each other prevail against anything evil can do. The Flowers and Gallants take it upon themselves to weave webs of trust, standing with friends and allies, obeying rulers, and guiding followers; in these webs they hope to catch the fitful gleams of Light that fall into the world and collect them into eternal beacons. Jewels seek positions of respect and authority in their community and exert their influence through the system. When they don't they build and tend new communities, write codes of law and establish customs for people who have none of their own, or have abandoned the societies they were born in. The ordered community, in which everyone has a place and no one is an exile, is the strongest possible bastion the Light can have.

\subsection*{Honour Tradition}

Tradition is a trust \textendash\, a store of wisdom laid in by past generations, to guide us in the present, and be handed on to the future. More than that, traditions are also our shared social identity. Shared traditions are what turns a group of people into a community, traditions builds trust between neighbours and a willingness to work for the communal good. The Flowers and Gallants do their best to breathe life into the rules and customs of the societies they find themselves in, especially the one they're born to. They are reluctant to flout a convention, especially when they don't know the reasons for it.  However, the Radiant are here to improve the world, not preserve it as it is, and when the Flowers see a Tradition they cannot condone they seek to replace it rather than remove it. The Gallants often use those Traditions they remember from past lives in the Kingdom when they need a substitute, but they can just as easily borrow something from the neighbouring culture (especially if they wish to bring them together), bring back a tradition from the \textquotedblleft good old days\textquotedblright, or formalise that awesome party their friends repeat every year. 

\section*{Duties}

A Flower's duties can be to build societies, to govern with grace and wisdom, to guard the traditions of the past, or even to stand vigilant at the border, but they always share one common theme. A Flowers duties revolve around people. 

Champions of Hearts lead from the front, they specialise in inspiring and leading their troops. On the battlefield tactical acumen and a strong commanding voice are as important as a good sword arm. Off the battlefield a Gallant must know how to keep moral high and foster a spirit of brotherhood among her troops. If she is skilled at her tasks a Gallant's fighting prowess shows itself through every one of her companions. 

Graces of Hearts are sometimes called Hearts of Hearts for the Queen and Calling fit together hand in glove. A Heart of Hearts builds societies, they lay down and exemplify the virtues people aspire too. They are both diplomats and leaders. It's a huge amount of work, but you are expected to have help from your friends. 

Menders of Hearts focus on human problems. They are more likely to be counsellors and therapists than doctors or nurses. The tools of a Mender are a good common sense understanding of people, words of encouragement and a winning smile. 

Seekers of Hearts seek to understand their societies from tip to toe. They can be historians unearthing long forgotten traditions, reporters and bloggers who speak about the lives of their fellows. Yet to the Gallants understanding is not for its own sake, it is for a purpose. They understand not only the unwritten rules of society, but the purpose of those rules and they use that understanding to purge corruption and fight abuses of the system in all their forms. 

Troubadours of Hearts incorporate the themes and virtues of their society into their work. They honour the heroes of old and create monuments to the ideals people aspire to. But at heart their work is simple, through their art the Troubadours bind people together with a strong group identity. 

\section*{Background}

Even before they Blossomed the followers of Hearts often had a lot of respect for the institutions of society. They (or at least their parents) hail from backgrounds in politics or law. Others cut their teeth with the codes and institutions of the playground or come from tight knit communities whose values fit neatly into the Court of Hearts. A lot of Princesses are attracted to the court because they wish to be one of the glamorous leaders they see among the Gallants but the Court's belief that leadership is hard work and duty means few stay. However the Court of Hearts does have a number of converts from the other Courts who felt working with Mortals and building societies provided a way to work on the big picture their own Court's had lacked. 

Regardless of their Background Flowers are always ok with the idea of others depending on them. Nearly all have some prior experience in being the responsible one, even if only to a younger brother or sister. While the Court doesn't exactly reject introverts or the socially unskilled it does expect them to work hard to overcome these limitations. Even the Gallients who have prior experience in taking responsibility for another have rarely held another's life in their hands, and they better start preparing for that first time. 

\section*{Signature Emotions: Trust and Duty}

Even the most powerful Princess with a blazing Inner Light cannot accomplish as much as a whole nation working together for their common good. Duty is the source of a Jewel's strength. By willingly subsuming herself into an ordered community where everyone knows their role and every role fits together like the stones of a magnificent castle, the Princess becomes something greater than herself. 

But there can be no duty without trust. A Jewel cannot obey orders from her superiors unless she trusts their judgement and their morality. She cannot delegate duties to her subordinates unless she trusts them to wield the authority she invests responsibly and competently. Duty is the stones, trust is the mortar that binds them. Separate they are worthless, but together they are what a better world is made of.

\section*{Character Creation}
Flowers and Gallants try to involve themselves in the lives of people around them, so many have above average Social Attributes. Those who actively strive to lead people emphasize Presence, while those who prefer more low-key or personal means focus on Manipulation. Empathy is a vital skill in order to notice the feelings, undercurrents, and motives of the people a Gallant wishes to help. Socialize and Persuasion are common, allowing a Flower to integrate herself into a group and sway others into doing the right thing and upholding the ideals of the Kingdom (whether they know of the Kingdom or not). Politics is also a common skill, used to improve the local political situation or simply use existing bureaucracies to their own ends. Subterfuge sees occasional use, mainly to keep important information from people who would use it selfishly or to protect the identities of a Princess's allies from those who would persecute them. 

Gallants who see merit in the society and political structure of the Kingdom frequently lean Academics to better understand both the society of the Kingdom and the modern world. Many old traditions generally involve art or music of one kind or another, so Expression is fairly common even for non- Troubadours. A Flower often favours Social then Mental and finally physical Attributes; Champions are the exception here, who emphasize their Physical attributes in order to fulfil their roles as protectors.

A Flower often places Social Merits over Mental Merits and Mental over Physical ones. Their ability to build and lead communities can lead to Circle, Mandate and a whole host of allies, contacts, retainers and staff. Most Flowers also master at least one conversation style, Noble Bearing is common and even Jewels who lead and cherish the traditions of anti-authoritarian or counter-culture groups have picked up some of the art through osmosis. Populist Rhetoric is also common but all styles have been seen in the Court of Hearts. The Small Unit Tactics merit is near universal for Champions of Hearts, as is merits such as Danger Sense and Inspiring to help keep the troops at maximum effectiveness. 

\subsection*{Heraldry}

Heraldry for the Nobles of Hearts strongly favours clothing in traditional styles, particularly formal dress. Elegant ball gowns and elaborate jewellery for phylacteries are very common among the Flowers, as are flowers or flower designs (roses especially) as accents, and pastels, rose-pink or -red, and shades of white. (It's not unknown for this sort of heraldry to be mistaken for a wedding dress.) Many Gallants follow the complementary mode of mens' formal wear, appearing in tuxedos, or in white tie, top hat and tails; or, of course, their own culture's equivilent. While modern formal is common a Princess of Hearts can adorn herself in the trappings of past civilisations to honour cultures she believes have valuable lessons for the present. The more practically minded (usually Champions) model their Heraldry on styles no less formal but less confining; the swashbuckling Cavalier, the wuxia hero, the knight in shining armour. These turn up on both sexes. In all cases, a Flower or Gallant is dressed well, showing off the best appearance of a high and civilized age.

\subsection*{Echo}
A Gallant's Echo makes it easier for people to get along, it helps smooth over many of those tiny unintentional bits of rudeness that can add up to ruining your day.
\subsection*{Practical Magic}
Through her leadership a Princess of Hearts inspires people to exceed their limitations. A Jewel has a pool of dice equal to the sum of her Social Attributes. She beings each scene with the  pool empty and may fill it once per scene by Reflexively spending a Wisp.

At any point the player may reflexively take dice from her pool and add them to rolls; however she cannot enhance her own rolls, only her ally's roll. She cannot add more than Terra dice to a single roll, and the recipient must be in the process of following the Princess' instructions. For this ability instructions must be kept short and simple: \textquotedblleft Investigate the mayor for corruption\textquotedblright\, is too broad, \textquotedblleft distract the mayor for a bit\textquotedblright\, or \textquotedblleft take a deep breath and tell me about the man who attacked you\textquotedblright\, is acceptable.

High Belief adds more dice to the Princess' pool. She gains two extra dice at Belief eight, nine and ten for a maximum of six extra dice. 
\subsection*{Invocation: Terra}
The principles of the Queen of Hearts find their magical expression in the Invocation of Terra, and her Flowers and Gallants learn it more easily than any other. It is bound up with the things the Queen values: mutual trust and peace, the responsible use of authority, sensitive compassion, and punctilious courtesies. It is also tied to earth and stone, especially when it has been refined and property constructed for human service, as proud cathedrals, elegant jewellery or solid stone walls.

Terra applies at no cost when the target of a Princess' Charm is earth or stone, including all forms of gemstones; and when her target is a non-supernatural human being who gives informed consent to the Charm. It also applies without cost when a Princess intends to resolve a conflict without violence, to make someone or something beautiful, to give requested aid without expecting any return, to coordinate the efforts of several people in a project, or to help a lawful authority in their mandated duties, or if you are a lawful authority undertaking your mandated duties.

A Princess of Hearts is a representative of her community and must always uphold its traditions and standards. Unless it would cause harm to another, breaking her community's social rules or standards of decorum blocks access to Terra for the rest of the scene

\section*{Quote}

We're not calling ourselves Princesses because our childhood fantasies came true. Nobility means something. 

\section*{Stereotypes}

\begin{itemize}
\item Clubs: Rough and unmannerly, but their hearts are in the right place.
\item Diamonds: Those towers of abstract reasoning leave me cold. You can't keep faith with real people by a theorem.
\item Spades: Yes, I suppose the mayor does deserve to have something happen to him ... why are you giggling?
\item Swords: It's easy to sacrifice yourself, and easier still when you know you'll be reborn. Living every day with duty, that's hard.
\item Tears: We both talk about nobility and duty. The difference is you use it as a justification, I remember what it actually means. 
\item Storms: At least the Swords only risk themselves ...
\item Mirrors: For the Light's sake, grow up and be responsible for once.

\item Vampires: Not all traditions are worth preserving.
\item Werewolves: Herd? Damn right I'm part of the herd. A herd billions strong who've claimed this whole earth. Now tell me, what is your little pack going to do about that?
\item Mages: You talk about Atlantis a lot but you only ever talk about its magic. What were the people like? Were the citizens content, the rulers just and the judges fair?
\item Prometheans: They share a culture, even when they haven't met another of their kind. There's something we're not seeing, something important.
\item Changelings: They had nothing to work with but scars and they built it into a society. Respect them for that if nothing else.
\item Sin-Eaters: You don't snub death by laughing in its face. You snub death by building something that will last long after you're gone.
\item Mummies: Our kingdom's have both fallen, but I'm trying to rebuild mine. 
\item Beasts: Teaching is a position of trust and authority. That is why we have so may rules about it, and you broke \textit{all} of them.
\item Mad Scientists: How did you convince him to turn himself in... An equation for the human mind? That can't be right, I don't believe it! I won't believe it!
\item Leviathans: Every city you build eats itself alive, you can't even live with your own family, and that's why I'm not afraid of you!

\item Hunters: If you claim to serve the Light then work in the light not the shadows.
\item Mortals: It's their world, they just need someone to teach them that.
\end{itemize}

\section*{Inspiration}
Queen Serenity, Minky Momo, Princess Celestia

\end{multicols}\subchapquote{The Queen of Spades}{
\noindent\textbf{AKA:} The Queen of Knaves, The Queen of Thieves\\
\noindent\textbf{Followers' Epithets:} Knaves (affectionately derogatory), Rogues (affectionately derogatory), Scoundrels (affectionately derogatory)\\
\noindent\textbf{Kingdom:} The Confederacy of Four Winds}
\phantomsection
\label{scha:Spades}\begin{multicols}{2}\raggedcolumns

\noindent \textit{Why does everyone think we are lazy? Being a Knave is hard work, last week I had five parties, five! That was a lot of planning, and dancing all night really tires you out. Then I spent all last night watching great stand ups and making notes. I know everyone in school's name, that took weeks of work, I had to write songs to help me remember, and then there's birthdays and everyone's favourite foods and what kind of jokes they like. After all that you have to keep yourself full of energy because any time something goes wrong people expect you to be the fast thinker or have a joke ready to soften the blow, and you can't do that when you're tired. }\\

\noindent \textit{But it's all worthwhile when I see my friends smile.}\\

\noindent A lot of the world's problems aren't nearly as hard as we \textquotedblleft know\textquotedblright\, they are. If we could just step outside the box and give ourself space to look at other ideas, space to breathe, we could accomplish so much. The Queen of Knaves leads a Confederation of Princesses who challenge common knowledge and teach people to laugh at the absurdities of the rules which define them. They do the Light's work with a quick wit, nimble hand and lateral thinking. 

\section*{Tales of the Kingdom}

Drifting above the mountain-side settlements that form The Confederacy of Four Winds the Queen of Spades rules from a palace carved out of the clouds themselves. Blessed with limitless energy the Queen breezes through her palace and amusements trailing courtiers and petitioners in her wake. Her regal presence comes not from her appearance but from her character. She is youthful, yet she possesses a presence that blows from her; carrying with it a sense of confidence and power, the smell of windswept desert sand, spices and the Queen's infectious laughter. Fashion at the royal court is a turbulent affair. The Queen changes her style frequently; always trying to keep one step ahead of fashions. This has memorably led to her occasionally preceding over a court of punks and ravers while dressed in ermine robes.

In conversation with her Princesses the Queen of Knaves breezes from topic to topic; she actively resists her Princess' trying to pin her down to a single topic. It's a game to her and unlike the citizens a Princess is expected to be skilled enough to play. Asking the Queen to stay focused (or worse: Get annoyed) is the fastest way to lose the game. If a Princess keeps her wits about her she may discover the Queen has answered all her questions (though not always in the order the questions were asked; and the answers may be hidden in riddles or metaphor). The Queen favours Princesses who play well, any Knave who can adopt her Queen's breezy style of casual conversation and trick the Queen of Knaves into answering a hidden question scores highly indeed.

Her critics say that the Queen of Knaves was a rebel before her coronation and remains a rebel at heart. Her supporters and citizens ask who better to lead a nation of tricksters, scoundrels and merchants than the Queen of Thieves herself? However the weight of truth travels with these accusations. The Queen is nostalgic for her days as a rebel and frequently leaves her throne unattended to travel her kingdom as a commoner, playing pranks and jokes before departing with the winds. She's infamously responsible for about half the graffiti in her kingdom. Her supporters again retort that what she lacks as an administrator, she makes up for by being the laughing heart of her people. 

\section*{Philosophy}

The Queen of Spades shows that you can solve any problem with an open mind and a smile.

\subsection*{Find the Centre of the Problem}

The Queen of Spades delights in unorthodox solutions. She has reached infamy for immediately suggesting one sentence solutions to complicated problems and insisting her followers prove it won't work or else carry it out. Many of these suggestions are obviously absurd but it doesn't matter, their purpose is not to solve the problem. By playing the devils advocate with a bad idea she forces her followers to question \textquotedblleft obvious\textquotedblright\, assumptions and ingrained dogma and so learn the first of her Philosophies: To look beyond the obvious and the assumed, to question what the real problem is, and to act only when she has found it. If the best answers are obvious after they've been found, then they are already obvious to she who just looks in the right way.

\subsection*{Free as the Wind}

A caged bird's music is never quite as sweet. The Queen of Spades says that Light is innate, let loose the chains and let it shine forth. A Rogue should always let her spirit soar and express herself fully, no matter what anyone might say and they should extend the same courtesy to everyone she meets. High above the clouds the Queen teaches her court how to live on their wits, their feet and beyond the reach of anyone who might chain them down. 

\subsection*{Laughter is the Best Medicine}

You can't change the past, but you can choose how you react to it. Who would prefer to cry than to laugh? To the Queen of Spades laughter really is the best medicine, and a positive attitude is just the thing to get a Princess who's feeling down out the door and fixing her problems. 

\section*{Duties}

Champions of Spades often rely on speed and agility in favour of raw power, if they have a weakness it's a tendency to rely on their ability to avoid a blow and consequently Knaves can't take a hit as well as some other Courts. When fighting alongside other Princesses many specialise in harassing their foes support while their allies fight the main battle. Outside of combat Champions of Spades are peerless scouts and messengers, before the spread of modern communication systems they served as a vital backbone in communication, trade and governance. Their powers of speed and flight made the knaves vital despite their carefree attitude.

Graces of Spades are not the teachers, mentors or diplomats you'd find in the other Courts. They don't want to change the way you think, a Grace of spades just wants to turn that frown upside down! The life of every party and the first onto the dance floor; everyone has a great time if there's a Grace of Spades around. Graces of Spades are often called upon by their counterparts in the other Courts, especially when they're working with large groups of mortals. It is easier to keep people's attention when they're enjoying being here but it's a double edged sword for the Knaves bring a party atmosphere that can be counterproductive when the Princess is working on serious matters.

Menders of Spades have an unusual focus, but a vital one. They heal the ills people inflict upon themselves. Closer to a teacher than a psychiatrist the Menders seek out those who've internalised harmful views of themselves; whether it's a cruel parent's insults or a demeaning cultural stereotype. They subtly and deftly they force their patents to confront their opinion again and again until they realise at last that it was no more solid than a dark cloud before the sun and vanishes in a puff of wind.

Seekers of Spades are taught to think laterally and to keep there eyes open. Often the answers to the toughest problems is just lieing out in the open, just waiting to be seen by someone open minded enough to accept it. In a Nakama a Seeker of Spades is the one to come up with the plan, directing the group to unorthodox but effective solutions, they do their best work with small focused goals. 

Whatever medium she uses, a Troubadour of Spades is the mistress of two arts: How to make people laugh and how to make people think. Some are out there day after day in the worst, most deprived areas showing people how to laugh despite it all. Others are renowned for art that looks at old ideas in new ways and finds fresh answers to questions that were long thought solved. The very best combine thought with laughter as peerless satirists who can have you laughing riotously at your ideological foes until you realise the jokes been on you. 

\section*{Background}

The free spirit of the Knaves is usually evident even before they Blossom. The Court of Spades draws its members from various countercultural groups (or as close as it can get given the young age of the average Blossoming), quite a few members were even outright criminals who Blossomed through a heroic effort to turn their life around.  Even when they aren't actually members of countercultural groups would-be spades often have some of the attitude. They are the self proclaimed jokers, slackers and class clowns who fight with authority, usually for the fun of it. 

Joining the Court of Clubs encourages these traits, but also softens them in its own way. Declaring war upon \textquotedblleft The Man\textquotedblright\, is less of an issue when you can just fly away, and when you can live on magic instead of a nine to five job. Knaves are more likely to make fun of what they see as wrong in society than they are to take up arms, and when they do decide to make a change their approach is often subtle, finessed and sometimes incomprehensible to others. It is usually delivered with a smile.

\section*{Signature Emotions: Laughter and Surprise}

In this tainted world of darkness the paths of least resistance and standard strategies are shrouded, to follow them will perpetuate the problems you wish to solve. Thus, say the Knaves, the best solutions will be surprising. They may be clear and logical after you've found them but at the moment of discovery you'll smack your forehead and ask how come you didn't think of this before. A true Knave must embrace surprise, both her own surprise at the moment of inspiration, and the surprise of those around her when she proposes unorthodox solutions, defies social conventions, or just lives freely in a world that's forgotten freedom. 

But Surprise is a fast and fleeting emotion. Something is needed to fill the gaps between each moment of inspiration and that something is laughter. A sense of humour will keep everyone's spirits up in this shrouded world, and laughing at herself takes away the sting of embarrassment whenever an unorthodox solution attracts ridicule or turns out to be unorthodox for a good reason. 

\section*{Character Creation}

As befitting a Court built in individualism and independence Knaves can be incredibly varied. Most have high scores in at least two Finesse attributes: Wits to live life outside the normal systems, Manipulation to get away with it. Dexterity isn't always favoured but if it is needed it tend to be high. Skills often emphasise the free spirited lifestyle: Survival to live on your own wits, Streetwise to make a fast buck without getting tied down to a nine to five job. Many knaves mix Persuasion and Subterfuge to talk their way out of trouble.

For Merits many Rogue's dashing charm and witty nature leads to a high Circle, which compensates for how her philosophy is at odds with the responsibilities of a Mandate. Knaves often have very broad Contacts, especially Graces for she who throws great parties shall never want for friends. Physically Fleet of Foot and Fast Reflexes come naturally to a practitioner of Aria. The Fast-Talking style is vital to get out of trouble in one peace.

\subsection*{Heraldry}

The Heraldry of Spades is the light blues of wind, the greys of the clouds and the oranges and browns of falling autumn leaves. The Knaves tend to be very individualistic in their choice of Regalia but light breezy outfits mixed with countercultural symbols are common. So are subtle visual puns for the observant viewer. 

\subsection*{Echo}

A knave's Echo tickles like a soft breeze, it's easier to laugh and easier to let go of what hurts you when a Princess of Spades is around. 

\subsection*{Practical Magic}

Knaves are unbound by any law except their own morality and have little problem evading the consequences imposed by those try to insist otherwise. Whenever the Knave is about to get into trouble she may Reflexively spend a Wisp to automatically avoid it. 

Trouble is defined broadly, she could avoid a teacher looking for an assignment she hasn't finished as easily as she could avoid a pack of roaming Darkspawn. Avoiding means that so long as the Princess focuses her full attention on the problem she will be unaffected for the remainder of the scene. So long as she continues to keep an eye out for the teacher he will not be able to find her but if she stops to gossip with friends he might sneak up on her. And of course nothing stops the teacher from waiting for the next scene and catching her in the classroom.

The methods Knaves use to stay out of trouble are varied, but this Practical Magic always favours the path of least magic and the path of least resistance (in that order). She might duck into a classroom with all the stealth of a professional spy or think up a lightning fast fib that sends the teacher away for a few moments while she makes a hasty exit but she wont avoid trouble by magically completing her essay in seconds or magically compel the teacher to let her off. 

When determining the path of least resistance this power takes into account the opposing character's mundane abilities but not their magical abilities, and if this results in the Knave trying a trick her opponent has a magical counter for a Clash of Wills ensues. \textit{Example: If that teacher was a Seeker of Diamonds who's Perception dice pool is superior to his dice pool for seeing through deception the Knave's magically enhanced instincts would direct her to spin a story only to trigger a Clash of Wills when the teacher used the Open Heart Charm.}

Once the Princess is in trouble, once the teacher begins dictating her punishment or the Darkspawn spot her and rolls Initiative for combat, then it becomes too late to use this ability to escape.

A Knave may invoke this power no more than Aria times per Session. High Belief increases this limit, granting one extra use a Belief eight, nine and ten.

\subsection*{Invocation: Aria}

The Invocation of Spades is Aria, the enviable lightness of being. It is the principle of detachment, the paradoxical unity of speed with forethought, and the laughter in the wind. Aria is taught by the Queen of Spades and her followers learn it easiest of all.

Aria applies at no cost when a Princess is blown by strong winds, is sprinting or moving at equivalent speeds under her own power, or is at risk of falling from a height; and when the target of her Charm is the air, or other gases. Aria also applies at no cost when the Princess aims to escape the consequences of her actions or help others do the same. It also applies without cost when a Princess intends to catch someone by surprise, to make someone laugh, to cast doubts, to expose pretensions and humiliate arrogance, and to subvert formal rules of an organisation or society.

Aria will not support any assertion of authority over another person. A Princess is free to advise, to persuade, to appeal to friendship, and even to deceive (within the limits of Belief) while using Aria, but if she issues a command she cannot apply the Invocation for the rest of the scene. %Add exemptions for life-death situations.

\section*{Quote}

Lighten up! This will be fine, and if not we'll just run for it.

\section*{Stereotypes}

\begin{itemize}

\item Clubs: Good-natured sort, by and large - except when they're not. The little groan from a “Let’s go clubbing” joke never gets old.
\item Diamonds: I probably shouldn't take quite so much pleasure in pointing out when they’re wrong. It’s a rare enough thing that it’s a treat, though.
\item Hearts: The bigger you make your castles the more places there are for me to hide in. You're far too easy a target my dear.
\item Swords: Halt dogooder! This is one of those bridges where you have to answer my riddles to pass.
\item Tears: You might have kept part of the Kingdom alive, but you haven't made it worth living. 
\item Storms: They are so predictable, traps, ambushes, they just charge straight in. 
\item Mirrors: Those special powers, did you know they stop working if... oops! Too late!

\item Vampires: They seem to be at the heart of any rotten power structure. I wonder; do they corrupt things, or are they just drawn to such things like flies to meat?
\item Werewolves: Discretion is often the better part of valour when a pissed-off eight foot monster is on the prowl. But sometimes you just have to get between them and their victims to draw them off.
\item Mages: Look at the puppet-masters, making society dance on strings to their tunes. Never let them catch you with scissors.
\item Prometheans: Well it's like... where to begin... I don't usually have to explain what a joke is.
\item Changelings: Ticksters and lawyers to a man. Don’t try to beat them at their game of words and promises unless you’re sure you’re very good, or they’ll ensnare you.
\item Sin-Eaters: Partying in the face of death; you've got style! Hey, I've never tried rum before. Want to share a bottle?
\item Mummies: What would happen if I woke you and commanded you to see the world and decide who you want to be? 
\item Beasts: Yeah, you're a troll, an internet troll. You're not scary, you're pathetic.
%\item Deviants: You know there's an easy answer. Start a new life in a Kingdom. Who needs a body or their memories right? You'll still be you and you can get the best treatment magic can provide. %This is practically a Xanatos gambit. If the Deviant agrees they escape their conspiracy and live a harmless happy life in the Dreamlands. If the disagree they realise they have something worthwhile to live for. 
\item Mad Scientists: Oh, don't mind us. We're just watching but um, can we like, be the beautiful lab assistants and shout \textquotedblleft It's Alive!\textquotedblright\, or something?
\item Leviathans: Um, right, yeah, living in the sky is starting to look like a good idea.
\item Hunters: Why are you so serious, ok I know why you're so serious, but you'd cope better if you laughed once in a while anyway.
\item Mortals: Come on everyone, smile! Fill my heart with sunshine!

\end{itemize}

\section*{Inspiration}
The Doctor, Psiren (Fullmetal Alchemist), Saint Tail, Captain Jack Sparrow, Ikuto Tsukiyomi, Pinkie Pie, V (V for Vendetta), Aladdin. 

\end{multicols}\subchapquote{The Queen of Swords}{
\noindent\textbf{AKA:} The Brightly Burning One, The Faithful Marshal, The Queen Errant\\
\noindent\textbf{Followers' Epithets:} Adventurers, Heroes, Martyrs\\
\noindent\textbf{Kingdom:} Aztallan}
\phantomsection
\label{scha:Swords}\begin{multicols}{2}\raggedcolumns

\noindent \textit{Fire, fire gets a bit of a bad name these days, but fire is really life. Almost every living thing on the planet owes its life to the sun, and we owe everything that makes life worth living to the fire inside us. Life is meaningless if we don't feel anything from it: Laughter, hope, wonder, even sadness. Without them what's the point?}

\textit{I guess, what I'm trying to say is. I love you, won't you come on an adventure with me? We'll see the world, try new things and maybe get into a few scrapes on the way. Even if we get burnt, pain is part of life and I wouldn't miss it for the world.}\\

\noindent Just as their magic lets them master the fire without, those who follow the Queen of Swords are consumed by the fire within. They are blazing heroic figures, driven by their passions. They will follow their hearts wherever it leads, even to their own deaths. The Court of Swords are called Martyrs too. 

\section*{Tales of the Kingdom}

In the jungles of Aztallan the Queen of Swords dwells in her capital built upon an artificial plateau atop a mighty pyramid. Here The Faithful Marshal rules from an obsidian throne and burns through the royal itinerary in bored and clipped, but supremely competent tones. She is a powerful looking woman, beneath her skin lies well-defined wiry muscles, her eyes positively smoulder with power. Outside the excitement of a crisis The Queen Errant makes little attempt to pretend she wouldn't rather be digging up new insects in the jungle; usually she appears before her court in rugged outdoor wear. The Queen of Swords is skilled at managing her image, she just usually doesn't care; but when it is required (and she always knows when it's required) she'll appear before her court dressed in anything from formal robes to full military uniform. Her household servants quickly follow suit, decorating the palace and statues of past heroes to match. Visiting Princesses are expected to dress and act with appropriate 
formalities until things return to normal.

The Faithful Marshal treats her Princesses with the same unpolished honesty she treats her court. Though she is never actually rude she won't hide her emotions if she feels the Princess' petition to be uninteresting or something that the Princess should have been able to solve on her own. If the Queen comes to befriend a Princess it colours her public interactions just as much as her private interactions. The Queen of Swords has the unique ability to simultaneously be running her kingdom and relaxing with her friends (or even canoodling on the throne).

Whenever the opportunity presents itself The Queen Errant leaves her throne and wanders her kingdom. Righting wrongs, doing odd jobs and exploring the jungles which she claims dominion over. She sleeps rough, or in whichever subject's house is close at hand. While her people are never quite comfortable with their Queen asking for a job, how does one give orders to his own Queen, they have at least had plenty of time to get used to her habits. When she can't get out the palace she enjoys study, gymnastics or spending time with one of her many consorts. It is said that long time residents of Aztallan can tell when the Queen falls in love anew from cities away by watching the wheels of government slowing down. Famously she will not kiss a man or woman unless she is deeply in love. 

\section*{Philosophy}

The Queen of Swords only ever gave one order. \textquotedblleft Whatever you do, do it out of love \textquotedblright. 

\subsection*{In the End, Only People Matter}

The Queen of Swords is not one for abstract morality and long Socratic debates about what is good and what is evil. Nor is she fond of formal systems of morality, duty or obligation, referring to them as either training wheels or a crutch. To the Queen of Swords morality can be reduced to a simple idea that everyone with Light in their hearts knows intuitively: If it hurts people then it is wrong. So don't do that. Furthermore The Queen believes that if one does not hold this truth in their heart then no amount of rules or regulations can make them good, but if a person is at heart good forcing them into an inappropriate social construct can crush their light. And so she teaches her followers to do what's right regardless of the rules, and she trusts them to know what's right. 

\subsection*{Specialisation is for Insects}

A Princess of Swords should be able to: Compose a sonnet, survive a shipwreck, design an aircraft, recite the epic poems, win her true love's heart, set a bone, clean a building, cook dinner, give orders, take orders, fight like a gentleman in the arena, fight dirty in the streets. The Philosophy is a lot more complicated than just saying you should be able to do everything, rather it consists of two related ideas. The first and simpler is simply to always be broadening your horizons, try new things, learn new skills. The second, like all The Brightly Burning One's Philosophies is about people. While she has nothing against relying on another, Queen of Swords teaches that dependence must always be a matter of convenience not necessity. True dependency is a chain, it allows others to force you away from being true to yourself. That path leads away from the light. 

\subsection*{Love Like a Flame}

The Brightly Burning One teaches her followers that their passions should blaze forth from them. They should feel with every fiber of their being: hope, joy, lust, sadness, fear, friendship, and especially love. Like a flame, their passions should light the way and bring warmth to the Princess and those around her. One person cannot save the world alone, but if she blazes bright enough she just might light the way. But a Martyr must always be aware that her love is like a flame and left unattended it will consume all before dying. The Princess should always know when it is time to let go lest she consumes herself. Yet when her very heart is at stake, her passions should burn everything in her path, even herself \textendash\, like a flame. 

\section*{Duties}

The followers of Swords are defined not so much by common duties but by a common approach to duties. They seek positions where they are free to follow their own methods and their own moral compasses rather than toe the line, and where their passion and energy can inspire those around them. Many see their duties as much about growing themselves as a person as it is about helping others and so seek to always find a new approach or a tougher challenge to stretch their limits.

Champions tend to find jobs outside of the normal social structures. They're more likely to be a private detective than a policewoman. More likely to specialise in deep sea rescue than be a fire fighter. They boast the most full time Darkspawn hunters of any Court. Champions of Swords are often trained to work independently, having the intelligence to plan, the strength to fight and the charisma to turn whoever is around them into an ally.

Graces often mentor on an individual level, seeking the maladjusted or the outcasts and helping them find where they can belong even if it's just a space within themselves. While often less versed in diplomacy than the other Courts Graces of Swords have some of the best connections to a variety of different cultures on a personal level thanks to the courts wanderlust and tendency to go adventuring. In more than a few cases they have turned someone's life around simply by encouraging them to move to a culture to which they are better suited. Graces who focus on larger scale projects are often exceptional at inspiring the public, but they usually ride of into the sunset once the immediate problem has been dealt with.

Menders are more likely to volunteer in a small African village than a major, regulated, hospital, those who work with machines often do the same working in oil rigs or digging wells, or they specialise in obscure but still important technology where their unique skills let them demand to work their way. In both cases they are likely to become experts at doing more with fewer resources.

Seekers often specialise in fields that blend academic learning with either demanding physical or social requirements. Archaeology in war torn parts of the world, private detectives knowledgeable in both forensic science and the seedy underworld, botanical trips through the Amazon rainforest or they are successors to the great and dangerous explorations of the past delving deep into the Underworld or the Shadow. More social exploits can include living with remote tribes for years or finding other cultures closer to home, no few Seekers of Swords have spent time living with dangerous criminals or supernatural predators simply to learn of their lives.

Troubadours of Swords are like a change on the wind. They breeze into town smelling of foreign and exotic ideas \textendash\, if you travel enough you can be foreign and exotic everywhere. Their art challenges norms but rarely opposes them (though some do take a delight in finding and shaking up repressed societies). She isn't saying the status qou is bad, only that it isn't the only way. Most take the effort to leave behind at least a few people who'll continue in the styles she teaches. 

\section*{Background}

Many members of Swords were adventurers before they Blossomed, they may not have faced terrible danger or performed heroic deeds but from backpacking across Europe to teaching English in China they showed the Court's spirit of adventure and self reliance. Others tapped into a confidence they never had before their Blossoming, when you gain magical powers and a new body that's stronger and smarter than your old one confidence is a natural result. Regardless of when they discovered it, Heroes are comfortable living on their own abilities and often eager for new experiences or just to push themselves further. 

Just as the Swords are the most adventurous of the Radiant, they are also the most passionate of the Radiant. Even before Blossoming they thought with their emotions, or really wished they had the confidence to do so. The court is overflowing with incurable romantics and people who are confident enough to follow their hearts into adventure, danger and whirlwind romances. While the Heroes are hardly stupid, they put little emphasis on personal safety when the alternative is an adventure. 

\section*{Signature Emotion: Love}

Love is not the oldest emotion, but it is the greatest of them all. Or so say the Court of Swords, for was it not with Love that nature created creatures that transcended base biological urges to survive and reproduce, creatures capable of selflessness and self sacrifice. What could be greater than that? And what better way to become a truly noble being capable of self sacrifice than finding something or someone you love?

\section*{Character Creation}

For Attributes Swords favour Intelligence, Strength and Presence equally and most are skilled at combining them in unusual ways to solve unusual problems. Heroes often have one or two dots in a wide variety of Skills to be self sufficent and because of their unusual, transient lifestyles. Survial is common among members who take the title of adventurer literally, as is Streatwise to find one's way around cities from London to Hong Kong. Many Heroes know how to fight, and Persuasion is often favoured for finding allies or peaceful solutions. 

For Merits the Martys often favour physical merits, social and even mental merits are harder to cultivate with an independent lifestyle. The Retainer Merit is surprisingly common, often representing a Sworn travelling companion. The philosophy of Swords makes it easy to form attachments. However it also teaches that one cannot live forever in another shadow so Princesses of Swords see this time as almost an apprenticeship where they teach  and broaden their companion's horizons with a grand adventure before parting ways.

\subsection*{Heraldry}

The court of Swords tends to favour blazing iconography. Reds, golds and whites are common. While it would not be accurate to call them flamboyant the Heroes wear their hearts on their sleeves. Their appearance often contains iconography that boasts of virtues they uphold and heroic deeds in their past and their Regalia radiates power. 

\subsection*{Echo}

A Hero's Echo inflames both passion and self confidence in those around her. She seems to light the path with her presence and blazing charisma. 

\subsection*{Practical Magic}

Love conquers all and a Princess of Swords can do the same. Once per scene she may spend a Wisp to automatically break through a barrier, whether she kicks down a door or finds words able to reach her love through the haze of madness. 
% For the purposes of this power love is defined broadly: Romantic love, familial love, friendship love, and courtly love all qualify, as do all other forms of love.

A barrier is defined as either a physical obstruction, a Condition, or a supernatural effect which prevents the pair from communicating properly. In turn proper communication is defined as the target's ability to comprehend the Princess, process what she is saying, and then act upon it.

If there is a seemingly mundane solution the Princess will try that first (or at least she should, otherwise she risks her secret identity). If she has appropriate mundane equipment it's equipment bonus increases by Fuoco. If the Princess leaps over a barrier she adds Fuoco feet to her maximum jumping distance. Otherwise she may try to tear through it and may do so if it's Durability is less or equal to (Strength – 4). If the barrier is supernatural or supernaturally reinforced a Clash of Wills ensures, if the Princess wins treat the barrier as mundane, otherwise she cannot breach it. Any damage inflicted is permanent.

Conditions or supernatural effects are merely suppressed for a short while. In time critical situations such as combat they are suppressed for Fuoco * 2 turns, in other situations Fuoco * 2 minutes. A Clash of Wills is required to deal with any supernatural influence.

High Belief makes the Princess' love shine brighter. She effectively increases her Strength and Fuoco by Belief - 7.

\subsection*{Invocation: Fuoco}


The Queen of Swords raises passionate devotion into magic in the Invocation of Fuoco, and the Adventurers learn it more easily than any other. Love, heroism and sacrifice are its fuel, the noblest things the Queen Errant knows. It is equally bound to literal fire, heat and light.

Fuoco applies at no cost when the target of a Princess' Charm is a flame, on fire, or hot enough to burn; and when the target is someone or something the Princess loves. It also applies without cost when a Princess intends to aid a person she loves, to strike at enemies with her full force, to take great risks for great gains, to inspire others to heroic endeavors, and at the culmination of a major campaign.


	\bigsidebarstart
	
	%\hspace*{-2cm}
    \centering \begin{tabulary}{1\textwidth}{|L|L|L|L|}
        \hline
        Nature                                               & Colour                & Strength             & Intensity        \\ \hline
        Love for an "idea" of a person not intimately known  & Translucent          & Acquaintance         & 1st degree (+0)  \\ \hline
        Romantic love                                        & red/orange           & Good friends         & 2nd degree (+1)  \\ \hline
        Parental love                                        & Purple                 & Lover, close family  & 3rd degree (+2)  \\ \hline
        Other familial love                                  & Blue                & Total devotion       & White hot (+3)   \\ \hline
		Courtly love                                         & Gold                & ~                    & ~ (+3)   \\ \hline
		Friendship love                                      & White                & ~                    & ~ (+3)   \\ \hline
        Obsession                                            & Shifting, unnatural  & ~                    & ~                \\ \hline
        \textquotedblleft As you wish\textquotedblright\,      & Rainbow              & ~                    & ~                \\
        \hline
    \end{tabulary}
	\bigsidebarend

The Invocation abandons Princesses who abandon their beloved. A Princess who betrays or otherwise injures a person she loves cannot apply Fuoco until the one she has hurt forgives her for the injury, or until a full lunar month has passed.

Several Fuoco Charms create flames out of the love one person feels for another. The nature of the love determines the colour of these flames, and the love's strength fuels their heat. Consult the table when a Princess uses these Charms. 

\section*{Quote}

I have no regrets, I'm doing this because I love you.

\section*{Stereotypes}

\begin{itemize}
\item Clubs: I think after spending so much time trying to be part of everyone else, they've forgotten who they are.
\item Diamonds: It's easy to say \textquotedblleft slow down and look at the big picture\textquotedblright\, when it's not your lover being held hostage.
\item Hearts: Sometimes traditions and etiquette are just a crutch. If you really know who you are and what mattered to you then you wouldn't need them. Not like they do.
\item Spades: Great for a laugh, I just wish they cared more deeply.
\item Tears: There's no love in what they do, it's nothing but fear.
\item Storms: Find something to love, then fight for it. If you do it the other way around you'll end like these girls. 
\item Mirrors: You're not supposed to be in love with yourself!

\item Vampires: I think the reason fire hurts them so much is that it reminds them what they've lost. They really are dead, on the inside.
\item Werewolves: If I was sure they wouldn't rage and splatter my intestines everywhere, I'd be their closest friend.
\item Mages: I don't care how long your theory is, it's still wrong. If anything is real it's what's in your heart, not magic.
\item Prometheans: See! I told you, it all comes down to the same thing. Fire and Humanity.
\item Changelings: You loved something so much it let you defeat gods and escape hell. You can't give up now!
\item Sin-Eaters: Keeping the flame alive in the face of death; I'd respect that, if I thought it meant anything to them.
\item Mummies: I'm only sort of immortal, but at least I'm alive. 
\item Beasts: \textquotedblleft If it hurts people then it is wrong.\textquotedblright\, That's all they do and all they are.
\item Mad Scientists: I think they are missing something. I think that they became what they are to try and find it, and they're still looking.
\item Leviathans: Everyone changes other people just by being near them. If we're all like beautiful stars orbiting each other then these are black holes drawing in and giving nothing back. Pure perfect evil.

\item Hunters: HEY! Why the hell are you shooting at me!? I'm on your side, here!
\item Mortals: I'm going to take you on an adventure and break you out of that shell. Just you watch me. 
\end{itemize}

\section*{Inspiration}
Sailor Uranus, Nanoha Takamachi, Iron Man, Indiana Jones, Robin Hood, Sayaka Miki, Aquaman (from the Batman the Brave and the Bold TV series)

\end{multicols}\subchapquote{Beacons}{}\phantomsection\label{scha:Beacons}\begin{multicols}{2}\raggedcolumns

\noindent \textit{Dear Diary}

\textit{I had the dream again. The one with me and Sir Rufus on a quest or the Heart of Earth that I always used to have when Alison bullied me. It made me feel much better, just like it always used to do.}

\textit{Maybe I had the dream because everyone's being so mean and pretending Jane doesn't exist. I should just talk to her, if I can make it through the Swamp of Sadness I can talk to one girl no matter what anyone else tells me!}\\

\noindent The Light of hope is far from unique to the Radiant; it belongs to all mankind. In the millennia of the Long Night mankind built civilisations and made wonders in the arts and sciences. It was regular humans who achieved the impossible and walked upon the moon, bringing a new dawn to end the Long Night. The Beacons are people whom the Light illuminates somewhat more brightly than most. Some Princesses say that in a better world everyone would be a Beacon. Some (usually the same people) say that in the Kingdom everyone was a Beacon.

A Beacon has a strong sense of hope, morality and idealism that radiates outwards. That’s all it takes - there are no astrological portents during their birth or inexplicable transformative events, only a person with a heart full of hope, empathy and kindness. Their presence encourages people to make the best of themselves. The effect is subtle enough that it could be taken as nothing more than the force of personality, if a Princess' magical sight did not show that Beacons carry a tiny spark of the Light.

Most Beacons don’t ever discover their supernatural potential. They live ordinary lives, and if they tend to be popular people surrounded by loyal friends or unusually sensitive to acts of cruelty that’s only what you'd expect to happen when someone is extraordinarily kind and emphatic. But Beacons do have their magic and their curses. Given the chance the Outer Darkness will go out of its way to kill a Beacon, and a Beacon carrying Shadows becomes vulnerable to the trap within the Dreamlands. Fortunately, as they entered the Dreamlands far from the Radiant Courts where the Wardens were strongest, many Beacons escaped the Wardens’ notice and instead discovered the strength to face problems in the waking world through their nocturnal experiences. 

Princesses usually keep an eye out for Beacons. If they seem in trouble they’ll offer help before the Beacon’s idealism fades and her Light is crushed. There are pragmatic reasons to keep an eye on Beacons: Their light tragically makes them a target to any Dark creatures they encounter, and they are disproportionately likely to Blossom. Being ready to ease a new Princess into her role can make a world of difference. But even if that weren’t so, the Hopeful would seek out Beacons anyway, just for the pleasure of their company.

\subsection*{The Nature of A Shining Heart}

A Beacons optimistic and empathic personality gives them an Echo that functions identically to a Nobles, Beacons use Integrity in place of Belief.

Like a Princess a Beacon may have Dreams in place of Aspirations, these Dreams do not need to fit a Calling since Beacons have no Calling. A Beacon may have up to three Aspirations or Dreams, in any combination, at a time. When a Beacon acquires Luminous Experience they may spend it on Integrity, removing Shadows, and Social Merits.

Beacons can use Bequests, they roll Integrity to transform a Bequest. As Beacons have no Wisps they are limited to Bequests with their own pool of Wisps or Bequests containing Charms that do not require Wisps to use. If the Charm's effects vary with Inner Light a Beacon is treated as having a single dot of Inner Light.

If a Beacon has the Girl Underground Condition (or the Dream Traveller merit) they may travel to the Dreamlands. They use the same mechanics as Princesses, though they lack the ability to protect themselves from the Gales with Transformation.

Because a Beacon is simply a regular person with a certain outlook on life Beacons qualify as mortal for all mechanical purposes. Supernatural powers that detect the presence of the supernatural will generally fail to notice a Beacon, unless they are especially sensitive or specialised to detect Light. Unlike an unTransformed Princess, a Beacon cannot hide their Light by storing inside their Transformed identity.

If a Beacon loses their sense of hope and idealism their Light fades and they become an ordinary mortal, they retain any Luminous Experience they've already acquired or spent but must replace any ongoing Dreams with Aspirations. Though they may keep the original Dream as an Aspiration. If they regain their hope they regain their Beacon status, as would any Mortal who became sufficiently kind, virtuous and idealistic (and buys the Virtuous Merit if they do not have it already). There are no hard rules for when a Beacon loses the template, but it will probably involve a Breaking Point. 

Finally Beacons suffer from Sensitivity, with a base of one die. Beacons use Integrity in place of Belief when calculating the effects of Sensitivity. 

\subsection*{Character Creation}

Beacons are given the Virtuous Merit for free, they may also buy the Sympathetic Merit for a single dot like Princesses. Beacons also have three free dots of Social Merits to represent the friends and connections they make by being so nice. In all other respects Beacons are created as ordinary mortals.

As an optional rule the Storyteller may allow protagonist Beacons to increase their three bonus merit dots to five dots and allow them to be spent on either social connections or supernatural abilities, be they innate abilities like Dream Traveller or Bequests. This rule is intended for Beacons who are aware of and make use of their supernatural nature, such that there are.

\begin{mybox}

\subsubsection*{Heir Apparent}

Reincarnated Princesses are born as Beacons before regaining their full powers. Most have the right attitude, and even if they don't the Beacon template is a suitable representation of a Princess' dormant power. Reincarnated Princesses have no special advantages over other Beacons, save that they cannot become any supernatural being that a full Princess could not.

\end{mybox}

\end{multicols}\subchapquote{Sworn}{}\phantomsection\label{scha:Sworn}\begin{multicols}{2}\raggedcolumns

\noindent \textit{Everyone thinks I should be jealous of my big sister because she's the one who can transform into a beautiful princess and has really cool magic.}

\textit{But they don't understand. They're not the ones who have to help her sew her wounds closed, they don't wake up early to help her cover up scars or hold her when she cries all night. I love my sister, she's amazing and works so hard for everybody, I wouldn't want to be her.}

\textit{But I will always be there for her.}\\

\noindent Wouldn’t it be nice if everyone was a Princess? If nobody could commit cruelty without suffering from Sensitivity? If encroachment of the Darkness could be beaten back by the magic of billions? If wishes were Blossomings then everyone would be a Princess already, but the Nobility know it’s no good wishing for a better world. You have to roll up your sleeves and get to work. Even if there is no way to force a Blossoming, a Princess with the Accept Fealty Charm can create a Sworn.

Even though Sworn lack the Inner Light of the Nobility they can still draw upon the Beliefs of their Queen like a Princess. A Sworn can feel his Queen's Invocation like a guiding light, it gives him a sense of hope and purpose that lets him push himself beyond human limits. That lets him perform magic. 

While a Queen's light is powerful enough to endow ordinary humans with magic it is also distant, shining from the far away Dreamlands. A Princess' Light shines from just a few feet away. For a Sworn working with a Princess they know and trust is an uplifting experience. Her Light shines upon the Sworn, filling him with hope and magic.

The Sworn fill many roles. Sometimes a Princess is impressed by a Mortal (or a Beacon) and offers them magic. This is an significant investment: not only must the Princess invest a portion of her power to create a Sworn, but she must make Bequests to provide the Sworn with Wisps and any Charms she wishes to offer. Other Sworn work closely with the Nobility. A team of Princesses who notice a gap in their skills could seek out a new member who has the mundane abilities. The powers of a Sworn and a few Bequests can provide the magical skills.

Some Princesses create Sworn for special tasks. Consider if you will a Princess who creates Bequests that grant intuitive insights into the working of legal codes and hands them out to a think tank drafting a complex new social justice law. If it works a simple investment could have an enormous long term payoff but such plans are risky. The Princess usually ends up running against one or more of the vested power groups within the World of Darkness.

Other Princesses just like to share their magic with people. Granting the powers of a Sworn to a beloved friend or family member can help another understand the Princess’ new life and preserve the bond between them. It is unfortunately not that uncommon for a friend or family member to be jealous and giving them some power can help smooth things over. A Princess who worries for her loved one’s safety can do worse than giving them the power to run away at superhuman speeds, or to work the magical defences around their shared home.

Among the Radiant Courts most Sworn got their position by sharing a close bond with a Princess. A friend, a family member, a rescued civilian who's gratitude blossomed into friendship and loyalty. Creating and supplying a Sworn is a big investment, and working closely with someone often means they'll learn your secret identity, so the Radiant have two good reasons to select someone they trust. Besides, when you think with your heart it often tells you to chose someone you care about. 

In Alhambra Sworn have served as a vital and prestigious pillar of society. The kingdom of Tears possesses vaults filled with Bequests, some dating back to the Kingdom itself. The Queen can produce Sworn in greater numbers than a Princess (but even her capabilities isn't infinite) and the reuse of old Bequests means the Court of Tears is able to maintain a large standing force of Sworn. Alhambra's civilisation depends on Sworn, Bequests and magic much as we depend on technology.

The armies of Storms also make much use of Sworn, who seem to form spontaneously from ordinary people who hold Tempesta’s fury in their hearts. These Sworn seem to be self sufficient, but lacking Bequests their powers are limited.

The rarest kind of Sworn is an deliberate attempt to redeem someone by using Sensitivity to show them  the truth of their actions. This is rare not only because creating a Sworn is a significant effort, but also because no one can become a Sworn without giving full and informed consent. Still people who believe their actions to be just may well consent to Sensitivity.

\subsection*{The Nature of a Loyal Companion}

Sworn draw their power from their Queen. When they act in accordance with their Queen's Philosophy, they can feel their magic grow stronger. And when they act against it, their powers weaken. Every Sworn follows a Queen and may purchase that Queen's Invocation at an out of affinity rate. Sworn also have access to practical magic based on their Court, but their powers are less potent than a full Princess' (and yes, the Court of Mirrors offers no useful Practical Magic to its Sworn).

Sworn can use Bequests, they roll Integrity to transform a Bequest. If the effects of the Charm within vary with Inner Light a Sworn has an effective Inner Light of one. 

Sworn may travel to the Dreamlands but (barring the Dream Traveller Merit) they require the Girl Underground Condition or a Princess to show them a Crawlspace entrance every time. Once they have seen a Crawlspace entrance they use the same mechanics as Princesses; though they lack the ability manifest Regalia. 

Sworn have a Wisp pool equal to Integrity. They may regain Wisps through the Dedication Merit, alternatively any Princess of the same Court may transfer Wisps through a ritual that re-confirms the Sworn's oaths. The ritual requires a minute to perform, and the full concentration of both parties.

A Sworn may leave the Queen's service at any time, becoming mortal once again. This may be a precursor to joining another Court. Experience spent on traits inapplicable for mortals are refunded.

For mechanical purposes a Sworn qualifies as a normal mortal under an enchantment. A supernatural power that detects supernatural beings will not register a Sworn, but one that detects the presence of active supernatural effects will. Powers that detect the Light will also detect a Sworn. 

Finally Sworn suffer from Sensitivity, with a base of one die. Sworn use Integrity in place of Belief to calculate the effects of Sensitivity.

\subsection*{Character Creation}

Sworn begin play with a single dot in their Queen's Invocation. In all other respects they are created as regular mortals. 

As an optional rule the Storyteller may allow protagonist Sworn to begin play with five bonus merit dots which may be spent on supernatural abilities, be they innate abilities like Dream Traveller or Bequests.


\begin{mybox}

\subsubsection*{Sworn Beacons}

Becoming a Sworn does nothing to damage a Beacon's optimistic personality, when a Beacon becomes Sworn they gain both sets of powers which makes Beacons a popular recruit for the upwardly mobile Princesses. 

Mechanically start with the Beacon template, then add a Sworn's Invocation, Wisp pool and practical magic. 

\end{mybox}


\subsection*{Practical Magic}

Sworn use a different set of Practical Magic than Princesses: 

\begin{itemize}
\item Clubs: The Queen of Clubs teaches her Sworn that Harmony comes from adapting to their surroundings. At any time they may spend a Wisp as an instant action to gain a Merit appropriate to their surroundings. He could acquire Danger Sense in a jungle or Scientific Rhetoric in a university but would be hard pressed to do the reverse. Merits acquired in this way cannot be supernatural in nature or represent an external asset like Allies or Resources. They last until the end of a Session a Sworn may only create up to Legno dots of Merits at any one time. 
\item Diamonds: The Queen of Diamonds cannot give the power of magic to all her people, so in it's place she grants the powers of technology and intellect. A Sworn of Diamonds can spend one Wisp to add his Acqua to a single roll to create equipment. This does not include Jury Rigging but it does include abstract equipment like plans and repositories.
\item Hearts: As a Princess of Hearts acquires positions of leadership and responsibility she is sure to want lieutenants that she can trust. A Sworn of Hearts may spend a Wisp to add the lower of their Status and Terra to any roll where they are acting to fulfil their duties within an organisation represented by the Status Merit. 
\item Spades: The Queen of Spades is an egalitarian above all else. Her Sworn are taught the same principles as her Nobles, within the limits of their magic. By spending a Wisp a Sworn of Spades may add Aria to one roll to create a false Clue or tamper with an existing Clue. They may also spend a Wisp to add Aria to an Athletics, Stealth, Streetwise or Subterfuge roll to get out of trouble but not on a roll to get into trouble. So they couldn't boost their Stealth to sneak into a secure building, but they could boost their Stealth to sneak back out again. 
\item Swords: Love knows no distance, and a Princess of Swords is always there to help her Sworn. By spending a Wisp a Sworn of Swords may add one third (round up) of the Transformed dice pool belonging to the Princess with whom she has the highest Dedication. She may only do this on rolls where she is the Primary Actor. This bonus is reduced by Intimacy, and though a Princess may have many loves she can still only be in so many places at once. Each trait may only assist one Sworn per turn. Naturally this power does not work if the Princess is already a Secondary Actor to the roll. \\ \textit{Example: A Sworn rolls Wits + Investigation and spends a wisp, adding one third of his Princess' Transformed Wits + Investigation as bonus dice. In the same turn another Sworn created by the same Princess spends a Wisp while rolling Wits + Empathy. he only adds a third of the Princess' Transformed Empathy. Both Sworn reduce the bonus dice by Intimacy to a minimum of -0.}
\item Tears:  As the Sworn of Alhambra protect their city they in turn are protected by the Queen of Tears. A Sworn of Tears can spend a Wisp to gain half Lacrima (round up) general armour for the remainder of the scene. 
\item Storms: The Queen of Storms' transformation lets her bestow greater gifts to her followers than any other Queen. This, and her laser focus means her Sworn share the same Practical Magic as Princesses of Storms. 
\item Mirrors: The True Heir has no need for assistance, and her sworn may cheer her on while she takes care of any problems. Consequentially the Sworn of Mirrors have a notably inferior Practical Magic: So long as they can see the Princess who swore them in a Sworn of Mirrors may spend a Wisp to add a single die to any roll made by that Princess.
\end{itemize}

\end{multicols}\subchapquote{Shikigami}{}\phantomsection\label{scha:Shikigami}\begin{multicols}{2}\raggedcolumns

\noindent \textit{“This symbol seems familiar. Do you think they took it from Andarta or did my memory of Andarta come from here?”}

\textit{“Maybe it's just a coincidence. It's a stag, lots of people put stags on their logos.”}

\textit{“Mmhmm. You can tell their only interest is money, the reception is marble and polished wood but the staff work in dark and cramped spaces. You should talk to their leader, it will be educational for you.”}

\textit{“Don't be silly, there's no way they'd let me in there.”}

\textit{“Of course they will, you're royalty. Making yourself belong anywhere is an easy Charm, but if you don't believe you have a right to talk to CEOs you'll never convince anyone else. Just walk in there like you know where you're going and have every right; otherwise I'll have to eat all those sweets you think you've hidden in that empty shoebox.”}\\


\noindent The Kingdoms may be fallen, the power of their mighty armies and enlightened Queens might only exist within the Dreamlands, but they still have their skills and experience.

Shikigami are Dreamlanders who have formed a symbiotic bond with a Princess. Her magic sustains the Shikigami allowing it to remain upon Earth indefinitely, in return it aids her with it's knowledge and magic reflected from her Light.

Most Shikigami were originally inhabitants of the Kingdom. They were men and woman of learning and renown who have pledged to aid the Nobility as they once did long ago; and yes, this does mean that in theory any living person could become a Shikigami but that almost never happens.

Those Shikigami who didn't come from the Kingdoms were regular Dreamlanders who were persuaded of the existence of the waking world, and chose to accompany a Princess on her way back. The real world may not add up like a mystery novel, but the astral reflection of Sherlock Holmes is still a great detective.

When the Release took place the Queens had plenty of time to prepare for the newly free souls to be born and grow until they were capable of Blossoming. Their agents scoured the Dreamlands for signs and portents and when they found omens of Nobles or future Nobles they permitted trusted advisor’s to possess animals so that they might find and offer tuition to young Princesses.

These days the loose Princess community shares the responsibility for finding freshly Blossomed Nobles and bringing them up to speed with Shikigami. Many Princesses are introduced to their Court by an experienced Princess or independent Shikigami and given a chance to meet potential candidates and choose her own full time advisor. As Nobles on Earth are able to meet their advisor in the Dreamlands there is rarely a need for Shikigami to inhabit animals. Today most Shikigami are inhabiting action figures, dolls, or plushies.

\subsection*{The Nature of a Cute Mentor}

A Shikigami contains the mind, but not the power, of an astral being woven into a physical shell by a tapestry of Light. While in this form they gain many of the powers a Princess possesses, just so long as they remain allied to a Princess who sustains them. Shikigami who possess an animal share their mind, as sentient and experienced beings their personality is dominant but many develop powerful but harmless vices such as excessive napping, a sweet tooth or chasing pieces of string. Even Shikigami who inhabit artificial animals such as teddy bears are affected by the astral resonance of their new body, usually to similar results. 

Despite their unusual body, a Shikigami is treated as a regular character with the full spread of Attributes and Skills. The Size modifiers (see page \pageref{SizeTable}) remain in permanent effect for a Shikigami. Unlike an animal, who's Attributes are designed to reflect a creature of their size, a Shikigami must pay extra and overcome the penalties if they want to pack human strength into a doll sized body.

All Shikigami have the dexterity of eight fingers and opposable thumbs in their mouth, paws, tail or whatever. If they are being observed Larcanay rolls are required to make this look like anything other than blatant magic.

Shikigami can swear themselves to a Queen. They may then learn her Invocation at an out of affinity rate and use Practical Magic as though they were a Sworn.

All Princesses have an innate ability to from a bond to a Shikigami, and much of a Shikigami's power depends upon this bond. 

Shikigami are skilled travellers of the Dreamlands, they use the same rules as a Princess except that whenever the Princess' Inner Light is added to a dice roll the Shikigami instead gains a flat +3 bonus. A Shikigami's Transformation protects against the gales, just like a Princess'. So long as a Shikigami is within the Dreamlands they may sever all ties to their bodies and Dedication to the Nobility and return to being a normal Dreamlander. A unique ability of Shikigami is to see and, if invited, enter the Crawlspace of any Princess they are bonded too. They can navigate the Crawlspace as a Teamwork action and when travelling to the Dreamlands either character can be the primary actor and use their own Intimacy to the destination. This makes Shikigami uniquely qualified to help a freshly Blossomed Princesses reach the Queens.

So long as it is bound to a Princes a Shikigami can Transform, and has access to Transformed Attributes and Skills. A Shikigami's maximum dots in a Transformed Attribute or Skill, and their maximum total Transformed dots, are both defined by the highest Inner Light among all the Princesses they are bonded too.

Shikigami have a Wisp pool equal to their Integrity. Shikigami automatically gain a Wisp every time a Princess they are bound too gains a Beat from an Aspiration or Dream (Dreams grant one only Wisp, not one for each kind of Beat). Princesses can an also transfer Wisps directly to their bonded Shikigami, this requires a physical contact and a minute. While touching is sufficient, petting, grooming or playing makes it easier. Without a Bond a Shikigami has no natural way to regain Wisps and must rely on somebody who knows the Charm Charge. 

While Transformed a Shikigami can use Charms, however the maximum Charm rating a Shikigami can learn is limited by the highest Inner Light among the Princesses a Shikigami is bonded to. At Inner Light one a Shikigami can learn one dot charms, every two Dots of Inner Light increases this limit by one. Shikigami have no affinity for any Charm trees. This pusedo Inner Light rating also applies when a Shikigami uses a Charm that's effects vary by Inner Light.

However the last and greatest ability of Shikigami is that they serve as magical relays for the Nobility. Any Princess bound to a Shikigami automatically has an Intimacy equivalent to Love to that Shikigami and every other Princess bound to her. Taken as a group they have a Commonality of Dedicated.

\subsection*{Character Creation}

Shikigami are created as an ordinary mortal character, except that their starting Size is somewhere between 1-3. Most are animate dolls, plushies or small animals. Shikigami can buy dots of Transformed Size, one point of Transformed Size costs one Merit Dot. A Shikigami can buy enough Transformed Size to reach Size eight, and no more.

If Shikigami is a player character he starts play with one free dot of White Rabbits, as astral beings Shikigami are intimately connected to the Dreamlands. NPC Shikigami experience prophetic dreams whenver the Storyteller wishes. 

Shikigami may buy physical abilities for their body as Merits. A Shikigami with wings can fly at it's usual speed for four merit dots, one with sharp teeth or claws may buy damage zero or one weapons for one or two Merit dots respectively. Physical abilities that emerge or enhance when a Shikigami Transforms should be represented by acquiring Charms. 

Finally Shikigami being play with five points of Experience which they may spend upon Merits (including physical abilities), Transformed Abilities or Skills, Charms or Invocations.

\subsection*{Binding a Shikigami}

Every Hopeful has the ability to form a bond with an inhabitant of the Dreamlands. To do this a Princess needs two things. The first is a Dreamlander who is willing to travel to Earth with her and who has a Rank of 5 or less (former inhabitants of the Kingdom don't have a Rank but they all qualify except the Queens). This is usually quite easy, the Dreaming Kingdoms have institutions whose soul purpose is to match Princesses with potential Shikigami. 

The second thing is a body for the Dreamlander to inhabit. Usually this means a doll or a plush toy. An animal can also be used but inanimate objects are favoured for ethical reasons. The body should resemble the Dreamlander, usually this means a thematic resonance, such as a toy fox for a shrewd and cunning Dreamlander. However if the Princess finds a Dreamlander who actually is a talking fox then naturally a toy fox will make an appropriate body regardless of his personality. 

When the Princses has both a recruit and a vessel for him she must go to the Dreamlands while her sleeping body cuddles the vessel. Then the Princess can perform the rite of binding upon the Dreamlander.

  \begin{myindentpar}{10pt}  Action: Extended, Inner Light + Presence (10 minutes/roll, threshold = 3x Dreamlander's Rank, 6 successes for former inhabitants of the Kingdom) 
    Cost: 1 Willpower 
    
    \emph{Dramatic Failure:} The Princess wakes without transforming the vessel. She can never make this Dreamlander into a Shikigami.
    
    \emph{Failure:} The Princess makes no progress. If she fails to bind the Shikigami she gains the Stuck Magic condition.
    
    \emph{Success:} The Princess makes progress. If she reaches the threshold the Dreamlander becomes a Shikigami, he takes on some of his new body's nature and in exchange gains the ability to travel between Earth and the Dreamlands at will. 
    
    \emph{Exceptional Success:} The Princess makes great progress.
	
    Suggested modifiers: the vessel is appropriate for the Dreamlander (+1 to +3), the vessel is inappropriate for the Dreamlander (-1 to -3) \end{myindentpar}

A Princess may only be bound to one Shikigami at a time; The Charm Astral Companion allows the Nobility to invite Dreamlanders to Earth she may still only share the Shikigami bond with one of her companions. 

If a Princess meets a Dreamlander who's visiting Earth with the aid of the Charm Astral Companion (or through some other means) she may use the above roll to create the full bond without needing to visit the Dreamlands. If the Dreamlander is already inhabiting an object or a small animal halve the required successes. 

\end{multicols}
